%=============================================================
%  REVISIÓN SISTEMÁTICA — Deep Learning para la detección de
%  anomalías en tráfico vehicular (sensores, trayectorias y video)
%  Guías: Kitchenham y Charters (2007); PRISMA 2020; PRISMA-S
%  Versión 2 (respuesta al reporte de revisión del 13/10/2026)
%
%  Las cifras de resultados se generan con
%      python3 ../suplementos/generar_conteos.py
%  y se cargan desde generado/*.tex. No escribir conteos a mano.
%  Los datos que solo pueden aportar los autores aparecen como
%  \pend{...} (en rojo en el PDF) y se listan en PENDIENTES.md.
%=============================================================

\documentclass[12pt, a4paper]{article}

\usepackage[utf8]{inputenc}
\usepackage[T1]{fontenc}
\usepackage[spanish, es-nodecimaldot, es-tabla]{babel}
\usepackage{lmodern}
\usepackage{microtype}
\usepackage{setspace}
\onehalfspacing
\usepackage[top=3cm, bottom=2.5cm, left=3cm, right=2.5cm]{geometry}
\usepackage{amsmath, amssymb}
\usepackage{graphicx, float, caption}
\captionsetup{font=small, labelfont=bf}
\usepackage{tikz}
\usetikzlibrary{arrows.meta, positioning, calc}
\usepackage{booktabs, tabularx, array, multirow, longtable}
\renewcommand{\arraystretch}{1.25}
\newcolumntype{Y}{>{\raggedright\arraybackslash}X}
\newcolumntype{P}[1]{>{\raggedright\arraybackslash}p{#1}}
\usepackage{pdflscape}
\usepackage{enumitem}
\setlist[itemize]{leftmargin=1.5em, itemsep=2pt}
\setlist[enumerate]{leftmargin=1.5em, itemsep=2pt}
\usepackage[table]{xcolor}
\definecolor{azulPrincipal}{RGB}{0, 70, 127}
\definecolor{azulMedio}{RGB}{0, 102, 179}
\definecolor{azulClaro}{RGB}{204, 229, 255}
\definecolor{grisClaro}{RGB}{245, 245, 245}
\definecolor{verdeOsc}{RGB}{0, 100, 60}
\definecolor{rojoAlerta}{RGB}{180, 0, 0}
\usepackage{tcolorbox}
\tcbuselibrary{skins, breakable}

\newtcolorbox{preguntabox}[1][]{enhanced,colback=azulClaro!50,colframe=azulPrincipal,
  fonttitle=\bfseries\color{white},coltitle=white,colbacktitle=azulPrincipal,
  title={Preguntas de investigaci\'{o}n},#1}
\newtcolorbox{hallazgo}[1][]{enhanced,breakable,colback=verdeOsc!8,colframe=verdeOsc,
  fonttitle=\bfseries\color{white},coltitle=white,colbacktitle=verdeOsc,
  title={Hallazgo},#1}

% \input primitivo: permite cargar filas dentro de tabular/longtable.
\makeatletter\newcommand{\tabinput}[1]{\@@input #1 }\makeatother

% Datos que deben completar los autores (aparecen en rojo).
\newcommand{\pend}[1]{\textcolor{rojoAlerta}{\textbf{[#1]}}}

\usepackage{fancyhdr}
\setlength{\headheight}{15pt}
\pagestyle{fancy}\fancyhf{}
\fancyhead[L]{\small\textcolor{azulPrincipal}{\textit{Deep Learning para la detecci\'{o}n de anomal\'{i}as en tr\'{a}fico vehicular}}}
\fancyhead[R]{\small\textcolor{azulPrincipal}{\textit{Art\'{i}culo de revisi\'{o}n}}}
\fancyfoot[C]{\small\thepage}
\renewcommand{\headrulewidth}{0.4pt}

\usepackage{titlesec}
\titleformat{\section}{\large\bfseries\color{azulPrincipal}}{\thesection.}{0.6em}{}[\titlerule]
\titleformat{\subsection}{\normalsize\bfseries\color{azulMedio}}{\thesubsection.}{0.5em}{}
\titleformat{\subsubsection}{\normalsize\itshape\color{azulPrincipal}}{\thesubsubsection.}{0.4em}{}

\usepackage{hyperref}
\hypersetup{colorlinks=true,linkcolor=azulPrincipal,citecolor=azulMedio,urlcolor=azulMedio,
  pdftitle={Deep Learning para la detecci\'{o}n de anomal\'{i}as en tr\'{a}fico vehicular: revisi\'{o}n sistem\'{a}tica},
  pdfauthor={W. A. Mayta Arotaype}}
\usepackage{cite}

%--- Cifras generadas desde la matriz de extracción ---------------------
% Generado por suplementos/generar_conteos.py. No editar a mano.
\newcommand{\Nestudios}{21}
\newcommand{\AnioMin}{2020}
\newcommand{\AnioMax}{2026}
\newcommand{\AnioA}{2}
\newcommand{\AnioB}{1}
\newcommand{\AnioC}{5}
\newcommand{\AnioD}{2}
\newcommand{\AnioE}{5}
\newcommand{\AnioF}{4}
\newcommand{\AnioG}{2}
\newcommand{\NrecientesVeintidosVeintiseis}{11}
\newcommand{\QAmedia}{4.38}
\newcommand{\QAde}{0.96}
\newcommand{\QAmin}{2.0}
\newcommand{\QAmax}{6.0}
\newcommand{\QAumbralAlta}{5.0}
\newcommand{\NaltaCalidad}{6}
\newcommand{\IDaltaCalidad}{S03, S04, S06, S10, S16, S17}
\newcommand{\QAUnosi}{6}
\newcommand{\QADossi}{7}
\newcommand{\QATressi}{12}
\newcommand{\QACuatrosi}{2}
\newcommand{\QACincosi}{2}
\newcommand{\QASeissi}{9}
\newcommand{\QASietesi}{11}
\newcommand{\QAOchosi}{6}
\newcommand{\Ncodigo}{13}
\newcommand{\NcodigoURL}{3}
\newcommand{\NetiqVerificadas}{6}
\newcommand{\NdatosSimuladosAny}{3}
\newcommand{\NdatasetPropio}{6}
\newcommand{\Nvariabilidad}{2}
\newcommand{\NdatoSensores}{4}
\newcommand{\NdatoMovilidad}{2}
\newcommand{\NdatoTrayectorias}{3}
\newcommand{\NdatoVideo}{7}
\newcommand{\NdatoMultimodal}{3}
\newcommand{\NdatoSimulado}{2}
\newcommand{\Nurbano}{7}
\newcommand{\Nautopista}{7}
\newcommand{\Nambos}{5}
\newcommand{\NentornoNE}{2}
\newcommand{\NfamRec}{14}
\newcommand{\NfamConv}{12}
\newcommand{\NfamAE}{6}
\newcommand{\NfamGrafo}{8}
\newcommand{\NfamTransf}{10}
\newcommand{\NfamGen}{2}
\newcommand{\Nhibrido}{17}
\newcommand{\Nmultimodal}{4}
\newcommand{\NetiqReal}{5}
\newcommand{\NetiqManual}{5}
\newcommand{\NetiqBench}{3}
\newcommand{\NetiqInyect}{4}
\newcommand{\NetiqMixto}{3}
\newcommand{\NetiqSin}{1}
\newcommand{\NetiqNR}{0}
\newcommand{\Ntasa}{9}
\newcommand{\Numbral}{19}
\newcommand{\Nparticion}{18}
\newcommand{\Nunidad}{20}
\newcommand{\Nmetaccuracy}{10}
\newcommand{\Nmetprecision}{7}
\newcommand{\Nmetrecall}{8}
\newcommand{\NmetFuno}{7}
\newcommand{\Nmetaucroc}{8}
\newcommand{\Nmetprauc}{2}
\newcommand{\Nmetdr}{4}
\newcommand{\Nmetfarfpr}{5}
\newcommand{\Nmetmttd}{2}
\newcommand{\Nmeterrormaermsemape}{4}
\newcommand{\Nmetioumap}{1}
\newcommand{\Nmetoperativaslatenciafpsenergia}{11}
\newcommand{\NcatA}{9}
\newcommand{\PcatA}{43~\%}
\newcommand{\NcatAnyA}{12}
\newcommand{\NcatB}{5}
\newcommand{\PcatB}{24~\%}
\newcommand{\NcatAnyB}{9}
\newcommand{\NcatC}{4}
\newcommand{\PcatC}{19~\%}
\newcommand{\NcatAnyC}{5}
\newcommand{\NcatD}{3}
\newcommand{\PcatD}{14~\%}
\newcommand{\NcatAnyD}{5}
\newcommand{\Nbaselines}{21}
\newcommand{\NbaselineNoDL}{11}
\newcommand{\NsinBaseline}{0}
\newcommand{\Nlimdatosreplicabilidad}{7}
\newcommand{\Nlimetiquetado}{9}
\newcommand{\Nlimgeneralizacion}{17}
\newcommand{\Nlimmetricasnouniformes}{1}
\newcommand{\Nlimcostodespliegue}{8}
\newcommand{\Nliminterpretabilidad}{1}
\newcommand{\XSensoresRec}{2}
\newcommand{\XSensoresConv}{1}
\newcommand{\XSensoresAE}{0}
\newcommand{\XSensoresGrafo}{2}
\newcommand{\XSensoresTransf}{2}
\newcommand{\XSensoresGen}{0}
\newcommand{\XMovilidadRec}{1}
\newcommand{\XMovilidadConv}{1}
\newcommand{\XMovilidadAE}{2}
\newcommand{\XMovilidadGrafo}{2}
\newcommand{\XMovilidadTransf}{0}
\newcommand{\XMovilidadGen}{0}
\newcommand{\XTrayectoriasRec}{3}
\newcommand{\XTrayectoriasConv}{2}
\newcommand{\XTrayectoriasAE}{1}
\newcommand{\XTrayectoriasGrafo}{2}
\newcommand{\XTrayectoriasTransf}{1}
\newcommand{\XTrayectoriasGen}{1}
\newcommand{\XVideoRec}{4}
\newcommand{\XVideoConv}{6}
\newcommand{\XVideoAE}{1}
\newcommand{\XVideoGrafo}{0}
\newcommand{\XVideoTransf}{2}
\newcommand{\XVideoGen}{0}
\newcommand{\XMultimodalRec}{3}
\newcommand{\XMultimodalConv}{2}
\newcommand{\XMultimodalAE}{1}
\newcommand{\XMultimodalGrafo}{1}
\newcommand{\XMultimodalTransf}{3}
\newcommand{\XMultimodalGen}{0}
\newcommand{\XSimuladoRec}{1}
\newcommand{\XSimuladoConv}{0}
\newcommand{\XSimuladoAE}{1}
\newcommand{\XSimuladoGrafo}{1}
\newcommand{\XSimuladoTransf}{2}
\newcommand{\XSimuladoGen}{1}
\newcommand{\NperiodoAntes}{3}
\newcommand{\XAntesRec}{1}
\newcommand{\XAntesConv}{3}
\newcommand{\XAntesAE}{0}
\newcommand{\XAntesGrafo}{1}
\newcommand{\XAntesTransf}{0}
\newcommand{\XAntesGen}{0}
\newcommand{\NperiodoMedio}{7}
\newcommand{\XMedioRec}{4}
\newcommand{\XMedioConv}{3}
\newcommand{\XMedioAE}{3}
\newcommand{\XMedioGrafo}{2}
\newcommand{\XMedioTransf}{2}
\newcommand{\XMedioGen}{1}
\newcommand{\NperiodoReciente}{11}
\newcommand{\XRecienteRec}{9}
\newcommand{\XRecienteConv}{6}
\newcommand{\XRecienteAE}{3}
\newcommand{\XRecienteGrafo}{5}
\newcommand{\XRecienteTransf}{8}
\newcommand{\XRecienteGen}{1}
\newcommand{\NgrafoNoVideo}{8}
\newcommand{\Nsupervisado}{10}
\newcommand{\Nnosupervisado}{11}
\newcommand{\NenfoquePred}{5}
\newcommand{\NenfoqueRecon}{5}
\newcommand{\NenfoqueClasif}{9}
\newcommand{\NetiqRealesTotal}{13}
\newcommand{\NetiqNoReales}{8}
\newcommand{\NumbralNoAplica}{8}
\newcommand{\NumbralExplicito}{11}
\newcommand{\NaltaEtiqReal}{3}
\newcommand{\NaltaHibrido}{5}
\newcommand{\NexcluidosLimite}{5}
\newcommand{\NincluidosLimite}{7}

% Generado por suplementos/prueba_sensibilidad_cadena.py. No editar a mano.
\newcommand{\NsensibOK}{19}
\newcommand{\IDsensibNo}{S03, S14}

%=============================================================
%  DATOS DE LA BÚSQUEDA Y DEL CRIBADO
%  Las fechas, los conteos por base, el flujo PRISMA y los kappa los
%  escribe automáticamente cribado/scripts/06_resultados_prisma.py.
%  A mano solo: \RegistroOSF y \RepositorioDatos.
%  Lo pendiente aparece en rojo en el PDF.
%=============================================================

% --- Fechas y registro ---------------------------------------
\newcommand{\FechaBusqueda}{19/10/2026}
\newcommand{\FechaBusquedaEN}{19 October 2026}
\newcommand{\RegistroOSF}{\pend{OSF, DOI 10.17605/OSF.IO/XXXXX}}
\newcommand{\RepositorioDatos}{\url{https://github.com/ClaudioB12/rsl-wilbert}}

% --- Revisor y fiabilidad del cribado (autor único) ------------
\newcommand{\Revisor}{W.\,A.\,M.\,A.}
\newcommand{\KappaRecribadoTA}{\pend{0,xx}}     % re-cribado ciego, título y resumen
\newcommand{\NRecribadoTA}{\pend{n}}
\newcommand{\KappaRecribadoTC}{\pend{0,xx}}     % re-cribado ciego, texto completo
\newcommand{\NRecribadoTC}{\pend{n}}
\newcommand{\KappaAsistenteTA}{0,06}     % acuerdo autor-asistente (TA)
\newcommand{\AcuerdoAsistenteTA}{31~\%}
\newcommand{\NExclRegla}{100}             % exclusiones propuestas por regla
\newcommand{\NExclReglaConfirmadas}{100}  % ... confirmadas por el autor
\newcommand{\KappaIATA}{0,77}            % acuerdo autor-IA en TA (pasa / no pasa)
\newcommand{\AcuerdoIATA}{93~\%}
\newcommand{\NCambiosIATA}{62}            % decisiones de la IA que el autor cambió
\newcommand{\NFueraFiltros}{123}           % quitados por filtros del protocolo aplicados por script
\newcommand{\NexclOA}{577}                 % excluidos en TA por no ser de acceso abierto (CI-2)
\newcommand{\FechaOA}{19/10/2026}             % fecha de verificación del acceso abierto
\newcommand{\AcuerdoCalidad}{\pend{xx~\%}}     % puntaje propuesto vs. verificado

% --- Registros por base (Tabla de bases) ----------------------
\newcommand{\NScopus}{498}
\newcommand{\NWoS}{193}
\newcommand{\NIEEE}{152}
\newcommand{\NScienceDirect}{337}
\newcommand{\NACM}{397}
\newcommand{\FScopus}{19/10/2026}
\newcommand{\FWoS}{19/10/2026}
\newcommand{\FIEEE}{19/10/2026}
\newcommand{\FScienceDirect}{19/10/2026}
\newcommand{\FACM}{19/10/2026}

% --- Flujo PRISMA 2020 ----------------------------------------
\newcommand{\Pidentificados}{1577}   % suma de las cinco bases
\newcommand{\Pduplicados}{337}
\newcommand{\Pcribados}{1240}
\newcommand{\PexclTA}{1137}
\newcommand{\Pbuscados}{103}        % solicitados a texto completo
\newcommand{\PnoRecuperados}{11}
\newcommand{\Pevaluados}{92}
\newcommand{\PexclTC}{20}
% Motivos de exclusión a texto completo (deben sumar \PexclTC)
\newcommand{\PexclCEdos}{2}      % CE-2 no es tráfico vehicular
\newcommand{\PexclCEtres}{7}     % CE-3 solo pronóstico/clasificación
\newcommand{\PexclCEcuatro}{2}   % CE-4 sin DL como método principal
\newcommand{\PexclCEcinco}{3}    % CE-5 tipo de documento / preprint
\newcommand{\PexclCIseis}{5}     % CI-6 sin evaluación cuantitativa de la detección
\newcommand{\PexclCEseis}{0}     % CE-6 información insuficiente
\newcommand{\PexclOtros}{1}      % CE-1, CI-2, CI-3
% Otros métodos (búsqueda v1 en ACM DL a texto completo)
\newcommand{\POtrosIdentificados}{0}
\newcommand{\POtrosIncluidos}{0}

% --- Base de recuperación de cada estudio (Anexo B) -----------
% La escribe cribado/scripts/06_resultados_prisma.py en generado/base_recuperacion.tex.
\newcommand{\BaseRecup}[1]{\ifcsname baserecup#1\endcsname\csname baserecup#1\endcsname\else\pend{base}\fi}

\IfFileExists{generado/base_recuperacion.tex}{% Generado por cribado/scripts/06_resultados_prisma.py. No editar a mano.
\expandafter\def\csname baserecupS01\endcsname{Scopus, Web of Science}
\expandafter\def\csname baserecupS02\endcsname{Scopus, Web of Science}
\expandafter\def\csname baserecupS03\endcsname{Scopus, Web of Science}
\expandafter\def\csname baserecupS04\endcsname{Scopus, Web of Science}
\expandafter\def\csname baserecupS05\endcsname{Scopus, Web of Science, IEEE Xplore}
\expandafter\def\csname baserecupS06\endcsname{ScienceDirect}
\expandafter\def\csname baserecupS07\endcsname{Web of Science}
\expandafter\def\csname baserecupS08\endcsname{ScienceDirect}
\expandafter\def\csname baserecupS09\endcsname{ScienceDirect}
\expandafter\def\csname baserecupS10\endcsname{Scopus, Web of Science, ScienceDirect}
\expandafter\def\csname baserecupS11\endcsname{Web of Science}
\expandafter\def\csname baserecupS12\endcsname{Scopus, Web of Science}
\expandafter\def\csname baserecupS13\endcsname{Scopus}
\expandafter\def\csname baserecupS14\endcsname{Scopus, ACM Digital Library}
\expandafter\def\csname baserecupS15\endcsname{Scopus, Web of Science}
\expandafter\def\csname baserecupS16\endcsname{ScienceDirect}
\expandafter\def\csname baserecupS17\endcsname{Scopus, Web of Science}
\expandafter\def\csname baserecupS18\endcsname{Scopus, ACM Digital Library}
\expandafter\def\csname baserecupS19\endcsname{Scopus, Web of Science}
\expandafter\def\csname baserecupS20\endcsname{Scopus, Web of Science}
\expandafter\def\csname baserecupS21\endcsname{Scopus, Web of Science}
\expandafter\def\csname baserecupS22\endcsname{Scopus}
\expandafter\def\csname baserecupS23\endcsname{Scopus}
\expandafter\def\csname baserecupS24\endcsname{Scopus}
\expandafter\def\csname baserecupS25\endcsname{Scopus, Web of Science}
\expandafter\def\csname baserecupS26\endcsname{Scopus, Web of Science}
\expandafter\def\csname baserecupS27\endcsname{Scopus, Web of Science, IEEE Xplore}
\expandafter\def\csname baserecupS28\endcsname{Scopus}
\expandafter\def\csname baserecupS29\endcsname{ScienceDirect}
\expandafter\def\csname baserecupS30\endcsname{Scopus, Web of Science}
\expandafter\def\csname baserecupS31\endcsname{Scopus}
\expandafter\def\csname baserecupS32\endcsname{Scopus, Web of Science}
\expandafter\def\csname baserecupS33\endcsname{ACM Digital Library}
\expandafter\def\csname baserecupS34\endcsname{ACM Digital Library}
\expandafter\def\csname baserecupS35\endcsname{ScienceDirect}
\expandafter\def\csname baserecupS36\endcsname{Scopus, Web of Science}
\expandafter\def\csname baserecupS37\endcsname{ScienceDirect}
\expandafter\def\csname baserecupS38\endcsname{Scopus, ScienceDirect}
\expandafter\def\csname baserecupS39\endcsname{ScienceDirect}
\expandafter\def\csname baserecupS40\endcsname{Scopus, Web of Science}
\expandafter\def\csname baserecupS41\endcsname{Scopus}
\expandafter\def\csname baserecupS42\endcsname{ScienceDirect}
\expandafter\def\csname baserecupS43\endcsname{Scopus}
\expandafter\def\csname baserecupS44\endcsname{Scopus, Web of Science}
\expandafter\def\csname baserecupS45\endcsname{Scopus, Web of Science}
\expandafter\def\csname baserecupS46\endcsname{ScienceDirect}
\expandafter\def\csname baserecupS47\endcsname{Scopus, Web of Science, ScienceDirect}
\expandafter\def\csname baserecupS48\endcsname{Scopus}
\expandafter\def\csname baserecupS49\endcsname{ACM Digital Library}
\expandafter\def\csname baserecupS50\endcsname{ScienceDirect}
\expandafter\def\csname baserecupS51\endcsname{Scopus, Web of Science, IEEE Xplore}
\expandafter\def\csname baserecupS52\endcsname{Web of Science}
\expandafter\def\csname baserecupS53\endcsname{ScienceDirect}
\expandafter\def\csname baserecupS54\endcsname{ACM Digital Library}
\expandafter\def\csname baserecupS55\endcsname{Scopus, Web of Science}
\expandafter\def\csname baserecupS56\endcsname{Scopus}
\expandafter\def\csname baserecupS57\endcsname{Scopus}
\expandafter\def\csname baserecupS58\endcsname{ScienceDirect}
\expandafter\def\csname baserecupS59\endcsname{Scopus, Web of Science}
\expandafter\def\csname baserecupS60\endcsname{Scopus}
\expandafter\def\csname baserecupS61\endcsname{Scopus, Web of Science}
\expandafter\def\csname baserecupS62\endcsname{Scopus}
\expandafter\def\csname baserecupS63\endcsname{Scopus, Web of Science, IEEE Xplore}
\expandafter\def\csname baserecupS64\endcsname{Scopus, Web of Science}
\expandafter\def\csname baserecupS65\endcsname{Scopus, Web of Science}
\expandafter\def\csname baserecupS66\endcsname{Scopus, Web of Science}
\expandafter\def\csname baserecupS67\endcsname{Scopus, Web of Science, ScienceDirect}
\expandafter\def\csname baserecupS68\endcsname{Scopus, Web of Science}
\expandafter\def\csname baserecupS69\endcsname{Scopus, Web of Science}
\expandafter\def\csname baserecupS70\endcsname{ScienceDirect}
\expandafter\def\csname baserecupS71\endcsname{Scopus, Web of Science}
\expandafter\def\csname baserecupS72\endcsname{Scopus, Web of Science, ScienceDirect}
}{}

%=============================================================
\begin{document}
%=============================================================

\begin{center}
  {\LARGE\bfseries\color{azulPrincipal}
   Detecci\'{o}n de anomal\'{i}as en tr\'{a}fico vehicular mediante \textit{Deep Learning} a partir de datos de sensores, trayectorias y video: una revisi\'{o}n sistem\'{a}tica de la literatura}\\[1em]

  {\large Wilbert Alex Mayta Arotaype}\\[0.6em]

  {\small
    Escuela Profesional de Ingenier\'{i}a de Sistemas,
    Universidad Peruana Uni\'{o}n, Juliaca, Per\'{u}.\\
    \href{mailto:wilbert.mayta@upeu.edu.pe}{wilbert.mayta@upeu.edu.pe}
  }\\[0.8em]
\end{center}
\noindent\rule{\linewidth}{1pt}

%--- RESUMEN --------------------------------------------------
\section*{Resumen}
\textbf{Antecedentes:} la detecci\'{o}n de anomal\'{i}as en tr\'{a}fico vehicular (accidentes, congesti\'{o}n no recurrente, conducci\'{o}n peligrosa o disrupciones) se aborda cada vez m\'{a}s con \textit{Deep Learning} (DL) sobre datos de sensores, trayectorias y video. Las revisiones previas son narrativas, se limitan a un tipo de dato o no cubren el periodo reciente.
\textbf{Objetivo:} sintetizar los modelos de DL empleados para detectar anomal\'{i}as en tr\'{a}fico vehicular urbano y de autopista, junto con las m\'{e}tricas y los protocolos de evaluaci\'{o}n, los tipos de anomal\'{i}a, los conjuntos de datos, las l\'{i}neas base de comparaci\'{o}n y las limitaciones reportadas.
\textbf{M\'{e}todo:} revisi\'{o}n sistem\'{a}tica seg\'{u}n Kitchenham y Charters, reportada conforme a PRISMA 2020 y PRISMA-S. Se buscaron estudios de 2020 a 2026 en Scopus, Web of Science, IEEE Xplore, ScienceDirect y ACM Digital Library (\'{u}ltima b\'{u}squeda: \FechaBusqueda). Un revisor seleccion\'{o} los estudios con el apoyo de un asistente de cribado basado en reglas y verific\'{o} su consistencia mediante un re-cribado ciego ($\kappa$ intra-revisor~=~\KappaRecribadoTA). La calidad metodol\'{o}gica se evalu\'{o} con un instrumento de ocho \'{i}tems adaptado a estudios de aprendizaje autom\'{a}tico (procedencia de etiquetas, partici\'{o}n sin fuga, l\'{i}neas base, variabilidad, disponibilidad de c\'{o}digo y datos, validaci\'{o}n externa, protocolo y limitaciones). Protocolo registrado en \RegistroOSF.
\textbf{Resultados:} se incluyeron \Nestudios{} estudios primarios (\AnioMin--\AnioMax). Predominan las arquitecturas recurrentes (\NfamRec{} estudios), convolucionales (\NfamConv{}) y de atenci\'{o}n o Transformer (\NfamTransf{}); \Nhibrido{} estudios combinan varias familias. La categor\'{i}a de anomal\'{i}a m\'{a}s estudiada es accidentes e incidentes (\NcatA{} estudios, \PcatA). Solo \NetiqVerificadas{} estudios eval\'{u}an con etiquetas reales verificadas; \NetiqNoReales{} usan anomal\'{i}as inyectadas, pseudo-etiquetas o ninguna etiqueta. Solo \Ntasa{} reportan la proporci\'{o}n de anomal\'{i}as, \Nvariabilidad{} reportan variabilidad entre ejecuciones y \NcodigoURL{} publican su c\'{o}digo. La puntuaci\'{o}n media de calidad fue \QAmedia{} sobre 8 (rango \QAmin--\QAmax).
\textbf{Conclusiones:} el tipo de dato condiciona la familia de arquitectura m\'{a}s que el tipo de anomal\'{i}a. La comparabilidad entre estudios est\'{a} limitada sobre todo por la procedencia heterog\'{e}nea de las etiquetas y por protocolos de evaluaci\'{o}n incompletos, m\'{a}s que por la elecci\'{o}n del modelo. Se requieren benchmarks con etiquetas verificadas y un reporte m\'{i}nimo del protocolo de evaluaci\'{o}n.

\noindent\textbf{Palabras clave:} detecci\'{o}n de anomal\'{i}as; tr\'{a}fico vehicular; \textit{Deep Learning}; sistemas inteligentes de transporte; revisi\'{o}n sistem\'{a}tica.

%--- ABSTRACT -------------------------------------------------
\section*{Abstract}
\textbf{Background:} anomaly detection in vehicular traffic (crashes, non-recurrent congestion, dangerous driving or disruptions) is increasingly addressed with Deep Learning (DL) on sensor, trajectory and video data. Prior reviews are narrative, limited to a single data type, or do not cover the recent period.
\textbf{Objective:} to synthesize the DL models used to detect anomalies in urban and highway traffic, together with the evaluation metrics and protocols, anomaly types, datasets, comparison baselines and reported limitations.
\textbf{Method:} systematic review following Kitchenham and Charters, reported according to PRISMA 2020 and PRISMA-S. Studies from 2020 to 2026 were searched in Scopus, Web of Science, IEEE Xplore, ScienceDirect and the ACM Digital Library (last search: \FechaBusquedaEN). A single reviewer selected the studies, supported by a rule-based screening assistant, and checked consistency through blinded re-screening (intra-rater $\kappa$~=~\KappaRecribadoTA). Methodological quality was assessed with an eight-item instrument adapted to machine learning studies (label provenance, leakage-free split, baselines, variability, code and data availability, external validation, protocol, and limitations). Protocol registered at \RegistroOSF.
\textbf{Results:} \Nestudios{} primary studies (\AnioMin--\AnioMax) were included. Recurrent (\NfamRec{} studies), convolutional (\NfamConv{}) and attention/Transformer (\NfamTransf{}) architectures predominate; \Nhibrido{} studies combine several families. The most studied anomaly category is crashes and incidents (\NcatA{} studies, \PcatA). Only \NetiqVerificadas{} studies evaluate on verified real labels; \NetiqNoReales{} use injected anomalies, pseudo-labels or no labels at all. Only \Ntasa{} report the anomaly ratio, \Nvariabilidad{} report run-to-run variability and \NcodigoURL{} release their code. The mean quality score was \QAmedia{} out of 8 (range \QAmin--\QAmax).
\textbf{Conclusions:} the data type shapes the choice of architecture more than the anomaly type does. Comparability across studies is limited mainly by heterogeneous label provenance and incomplete evaluation protocols rather than by model choice. Benchmarks with verified labels and minimum reporting of the evaluation protocol are needed.

\noindent\textbf{Keywords:} anomaly detection; vehicular traffic; Deep Learning; intelligent transportation systems; systematic review.

%=============================================================
\section{Introducci\'{o}n}
\label{sec:intro}
%=============================================================

El crecimiento de las ciudades y del parque automotor ha aumentado la complejidad de la gesti\'{o}n del tr\'{a}fico vehicular, tanto en redes urbanas como en autopistas. El flujo vehicular presenta alta variabilidad temporal, ruido de medici\'{o}n y dependencia espacial entre v\'{i}as, sensores e intersecciones, por lo que debe analizarse como un fen\'{o}meno temporal y espacio-temporal y no como una secuencia aislada de mediciones~\cite{djenouri2019survey,zhang2020urban}.

Los Sistemas Inteligentes de Transporte (\textit{Intelligent Transportation Systems}, ITS) integran detectores de lazo y de microondas, c\'{a}maras de vigilancia, dispositivos GPS y plataformas de monitoreo. Sus datos se representan como series temporales (flujo, velocidad u ocupaci\'{o}n en un punto), como datos espacio-temporales (varias ubicaciones conectadas por la red vial), como trayectorias o como secuencias de video~\cite{djenouri2019survey,santhosh2020anomaly}. En este contexto, la \textit{detecci\'{o}n de anomal\'{i}as} busca identificar desviaciones relevantes respecto al comportamiento esperado: accidentes, congesti\'{o}n no recurrente, cierres de v\'{i}as, conducci\'{o}n peligrosa o disrupciones causadas por eventos extremos. Detectarlas a tiempo mejora la respuesta a incidentes y la seguridad vial~\cite{elsahly2022systematic,razi2023deep}.

Tradicionalmente, este problema se abord\'{o} con algoritmos comparativos y de umbral, modelos autorregresivos y aprendizaje autom\'{a}tico convencional~\cite{elsahly2022systematic,olugbade2022review}. Estos enfoques encuentran dificultades con el desbalance entre eventos normales y an\'{o}malos, con la dependencia espacial y con los datos incompletos o ruidosos. El \textit{Deep Learning} (DL) ofrece arquitecturas capaces de aprender representaciones temporales (redes recurrentes), espaciales (redes de grafos), visuales (redes convolucionales) o de reconstrucci\'{o}n (autoencoders), y su uso para detectar anomal\'{i}as en series temporales se ha generalizado~\cite{darban2024deep}.

\subsection{Revisiones previas y vac\'{i}o que aborda esta revisi\'{o}n}
\label{sec:previas}

Existen revisiones previas sobre el tema, pero ninguna cubre a la vez el periodo reciente, los distintos tipos de dato y el an\'{a}lisis del protocolo de evaluaci\'{o}n (Tabla~\ref{tab:previas}). Los \textit{surveys} de Djenouri et al.~\cite{djenouri2019survey} y Zhang et al.~\cite{zhang2020urban} cubren la literatura anterior a 2020, con predominio de m\'{e}todos estad\'{i}sticos y de miner\'{i}a de datos. Santhosh et al.~\cite{santhosh2020anomaly}, Fang et al.~\cite{fang2024vision} y Lin et al.~\cite{lin2025deep} se limitan al video. Belhadi et al.~\cite{belhadi2020trajectory} se limitan a trayectorias. Darban et al.~\cite{darban2024deep} tratan series temporales gen\'{e}ricas, no espec\'{i}ficas de tr\'{a}fico. La \'{u}nica revisi\'{o}n sistem\'{a}tica identificada, la de ElSahly y Abdelfatah~\cite{elsahly2022systematic}, trata la detecci\'{o}n autom\'{a}tica de incidentes en general, no se centra en DL y no cubre el periodo 2022--2026. Ninguna de ellas examina c\'{o}mo se obtienen las etiquetas de anomal\'{i}a ni c\'{o}mo se fijan el umbral y la unidad de evaluaci\'{o}n, aspectos que en detecci\'{o}n de anomal\'{i}as pueden cambiar los resultados m\'{a}s que la propia arquitectura~\cite{darban2024deep}.

% Generado por suplementos/generar_conteos.py. No editar a mano.
\begingroup
\small
\begin{xltabular}{\linewidth}{lcp{2.2cm}p{2.6cm}Y}
\caption{Revisiones previas sobre detecci\'{o}n de anomal\'{i}as de tr\'{a}fico y su diferencia con esta revisi\'{o}n.}\label{tab:previas}\\
\toprule
\textbf{Revisi\'{o}n} & \textbf{A\~{n}o} & \textbf{Tipo} & \textbf{Datos cubiertos} & \textbf{Vac\'{i}o respecto a esta RSL} \\
\midrule\endfirsthead
\toprule
\textbf{Revisi\'{o}n} & \textbf{A\~{n}o} & \textbf{Tipo} & \textbf{Datos cubiertos} & \textbf{Vac\'{i}o respecto a esta RSL} \\
\midrule\endhead
\bottomrule\endfoot
Djenouri et al.~\cite{djenouri2019survey} & 2019 & Survey narrativo & Sensores y trayectorias & Literatura previa a 2019; escaso DL; sin video; no analiza etiquetas ni protocolo \\
Santhosh et al.~\cite{santhosh2020anomaly} & 2020 & Survey narrativo & Video de vigilancia & Solo video; literatura hasta 2019; no sistemática \\
Belhadi et al.~\cite{belhadi2020trajectory} & 2020 & Survey con benchmark & Trayectorias & Solo trayectorias, no exclusivo de tráfico; predominan métodos no profundos \\
Zhang et al.~\cite{zhang2020urban} & 2020 & Survey narrativo & Datos urbanos multifuente & Anomalías urbanas en general; no cubre 2021--2026 \\
ElSahly y Abdelfatah~\cite{elsahly2022systematic} & 2022 & Sistemática (PRISMA) & Detectores de flujo & Detección de incidentes en general, no centrada en DL; sin video ni trayectorias; no cubre 2022--2026 \\
Olugbade et al.~\cite{olugbade2022review} & 2022 & Survey narrativo & Incidentes (IA/ML general) & Enfoque amplio en IA/ML; no sistemática; no cubre 2023--2026 \\
Razi et al.~\cite{razi2023deep} & 2023 & Survey narrativo & Video (seguridad vial) & Centrado en video; no trata series de sensores; no sistemática \\
Fang et al.~\cite{fang2024vision} & 2024 & Survey narrativo & Video (accidentes) & Solo video y solo accidentes; no sistemática \\
Zamanzadeh Darban et al.~\cite{darban2024deep} & 2024 & Survey narrativo & Series temporales genéricas & No específico de tráfico; sin video ni trayectorias \\
Lin et al.~\cite{lin2025deep} & 2025 & Survey (preprint) & Video (anticipación) & Solo anticipación de accidentes en video; preprint \\
Alinezhad Noghre et al.~\cite{alinezhadnoghre2025survey} & 2025 & Survey (preprint) & Video (humano, vehículo, entorno) & No exclusivo de tráfico; solo video; preprint \\
\end{xltabular}
\vspace{-6pt}{\footnotesize\raggedright Revisiones identificadas mediante b\'{u}squeda dirigida (2018--2026). Los \textit{preprints} se indican como tales.\par}
\endgroup


Esta revisi\'{o}n se diferencia de las anteriores en tres aspectos: (i) es sistem\'{a}tica y reproducible (PRISMA 2020 y PRISMA-S) y cubre 2020--2026; (ii) integra en un mismo an\'{a}lisis los estudios basados en sensores, trayectorias, movilidad agregada y video, lo que permite relacionar el tipo de dato con la familia de arquitectura; y (iii) extrae y sintetiza, para cada estudio, la procedencia de las etiquetas, el protocolo de evaluaci\'{o}n y las l\'{i}neas base, y eval\'{u}a la calidad metodol\'{o}gica con un instrumento espec\'{i}fico para estudios de aprendizaje autom\'{a}tico.

Las contribuciones son: (1) una matriz de extracci\'{o}n abierta de los estudios incluidos, de la que se derivan con un script todas las cifras de este art\'{i}culo; (2) un an\'{a}lisis cruzado entre tipo de dato, familia de arquitectura, tipo de anomal\'{i}a y periodo; (3) una s\'{i}ntesis del protocolo de evaluaci\'{o}n (procedencia de etiquetas, umbral, unidad de evaluaci\'{o}n y partici\'{o}n) y de las l\'{i}neas base; y (4) una agenda de investigaci\'{o}n basada en esas brechas.

\begin{preguntabox}
  \begin{enumerate}[label=\textbf{RQ\arabic*:}, leftmargin=3em]
    \item \textquestiondown Qu\'{e} t\'{e}cnicas de \textit{Deep Learning} se han utilizado para detectar anomal\'{i}as en tr\'{a}fico vehicular urbano y de autopista, y c\'{o}mo se relacionan con el tipo de dato?
    \item \textquestiondown Qu\'{e} m\'{e}tricas y qu\'{e} protocolos de evaluaci\'{o}n (procedencia de etiquetas, umbral, unidad de evaluaci\'{o}n y partici\'{o}n) se emplean para medir el rendimiento de estos modelos?
    \item \textquestiondown Qu\'{e} tipos de anomal\'{i}as de tr\'{a}fico se han abordado?
    \item \textquestiondown Qu\'{e} conjuntos de datos se han utilizado y con qu\'{e} tipo de etiquetas?
    \item \textquestiondown Qu\'{e} limitaciones y desaf\'{i}os reportan los estudios?
    \item \textquestiondown Contra qu\'{e} l\'{i}neas base se comparan los modelos de \textit{Deep Learning} y qu\'{e} ventaja reportan?
  \end{enumerate}
\end{preguntabox}

El resto del art\'{i}culo se organiza as\'{i}: la Secci\'{o}n~\ref{sec:marco} presenta el marco conceptual; la Secci\'{o}n~\ref{sec:metodo}, el m\'{e}todo; la Secci\'{o}n~\ref{sec:resultados}, los resultados por pregunta; la Secci\'{o}n~\ref{sec:discusion}, la discusi\'{o}n y las amenazas a la validez; y la Secci\'{o}n~\ref{sec:conclusiones}, las conclusiones.

%=============================================================
\section{Marco conceptual}
\label{sec:marco}
%=============================================================

\subsection{Datos de tr\'{a}fico y su representaci\'{o}n}

Una \textit{serie temporal de tr\'{a}fico} es una secuencia de observaciones de una variable de movilidad (flujo, velocidad u ocupaci\'{o}n) registrada por un detector en un punto de la red; puede ser univariada o multivariada. Cuando se observan simult\'{a}neamente varias ubicaciones, los datos son \textit{espacio-temporales} y suelen representarse como un grafo cuyos nodos son sensores, segmentos o zonas, y cuyas aristas son conexiones viales o relaciones origen-destino~\cite{djenouri2019survey}. Las \textit{trayectorias} son secuencias de posiciones de un veh\'{i}culo obtenidas por GPS o por seguimiento en video~\cite{belhadi2020trajectory}. El \textit{video de vigilancia} es una secuencia de im\'{a}genes de c\'{a}maras fijas o a bordo~\cite{santhosh2020anomaly}. El tipo de dato condiciona qu\'{e} anomal\'{i}as son observables y qu\'{e} arquitectura es adecuada.

\subsection{Anomal\'{i}as de tr\'{a}fico}

Una \textit{anomal\'{i}a de tr\'{a}fico} es un evento o patr\'{o}n que se aparta del comportamiento esperado de la red o de un veh\'{i}culo. En esta revisi\'{o}n se agrupan en cuatro categor\'{i}as: (i) accidentes e incidentes; (ii) congesti\'{o}n o flujo an\'{o}malo no recurrente; (iii) trayectorias o conducci\'{o}n peligrosa (contramano, exceso de velocidad, maniobras de riesgo); y (iv) disrupciones y eventos extremos (cierres de v\'{i}as, eventos meteorol\'{o}gicos, fallos de infraestructura). Las anomal\'{i}as son poco frecuentes, lo que produce conjuntos de datos desbalanceados y dificulta tanto el entrenamiento como la evaluaci\'{o}n~\cite{darban2024deep}.

\subsection{Enfoques de \textit{Deep Learning} para la detecci\'{o}n de anomal\'{i}as}

Los modelos de DL detectan anomal\'{i}as mediante cuatro mecanismos, que no dependen de la arquitectura concreta~\cite{darban2024deep}:
\begin{itemize}
  \item \textbf{Predicci\'{o}n:} el modelo (p.~ej., LSTM, GRU, TCN o Transformer) pronostica el comportamiento esperado y marca como an\'{o}malas las desviaciones grandes entre lo predicho y lo observado.
  \item \textbf{Reconstrucci\'{o}n:} un autoencoder, posiblemente de grafos o variacional, aprende a reconstruir datos normales y usa el error de reconstrucci\'{o}n como puntuaci\'{o}n de anomal\'{i}a. Es \'{u}til cuando hay pocas etiquetas.
  \item \textbf{Clasificaci\'{o}n supervisada:} con etiquetas de eventos normales y an\'{o}malos, una red (CNN, LSTM, GNN o h\'{i}brida) aprende directamente la frontera entre clases. Es frecuente en video y en detecci\'{o}n de incidentes con registros hist\'{o}ricos.
  \item \textbf{Similitud o contraste:} el modelo compara representaciones (por ejemplo, de video y de texto) y decide por similitud.
\end{itemize}
Los tres primeros mecanismos requieren un \textit{umbral} o una regla de decisi\'{o}n para convertir la puntuaci\'{o}n en una alarma. La elecci\'{o}n del umbral, junto con la \textit{unidad de evaluaci\'{o}n} (punto, ventana, frame, evento o trayectoria), determina las m\'{e}tricas resultantes; por eso esta revisi\'{o}n los extrae de cada estudio.

%=============================================================
\section{M\'{e}todo}
\label{sec:metodo}
%=============================================================

La revisi\'{o}n sigue las directrices de Kitchenham y Charters~\cite{Kitchenham2007} y se reporta conforme a PRISMA 2020~\cite{page2021prisma} y a su extensi\'{o}n para la b\'{u}squeda, PRISMA-S~\cite{rethlefsen2021prismas}. La lista de verificaci\'{o}n PRISMA 2020 con la ubicaci\'{o}n de cada \'{i}tem se presenta en el Anexo~\ref{ap:prisma}.

\subsection{Protocolo y registro}

El protocolo (PICOC, preguntas, criterios, estrategia de b\'{u}squeda, instrumento de calidad y plan de s\'{i}ntesis) se gestion\'{o} en Parsif.al y se registr\'{o} en Open Science Framework (\RegistroOSF). Las desviaciones respecto de la versi\'{o}n inicial (ampliaci\'{o}n del alcance a video y autopistas, sustituci\'{o}n de MDPI y Taylor \& Francis por IEEE Xplore, regla general sobre \textit{preprints}, nuevo instrumento de calidad y nueva pregunta RQ6) se documentan en el historial del protocolo (Anexo~\ref{ap:protocolo}, Tabla~\ref{tab:historial}).

\subsection{PICOC y alineamiento del alcance}

\begin{table}[H]
\centering
\small
\caption{Marco PICOC de la revisi\'{o}n.}
\label{tab:picoc}
\begin{tabularx}{\linewidth}{>{\bfseries}lY}
\toprule
Elemento & Descripci\'{o}n \\
\midrule
Poblaci\'{o}n    & Tr\'{a}fico vehicular en redes urbanas y autopistas, observado mediante sensores fijos, trayectorias GPS, datos de movilidad agregada o video de vigilancia \\
Intervenci\'{o}n & Modelos de \textit{Deep Learning} (recurrentes, convolucionales, autoencoders, grafos, Transformers, generativos o h\'{i}bridos) usados como m\'{e}todo principal para detectar anomal\'{i}as \\
Comparaci\'{o}n  & L\'{i}neas base reportadas por cada estudio: otros modelos de DL, aprendizaje autom\'{a}tico convencional, m\'{e}todos estad\'{i}sticos o reglas (RQ6) \\
Resultado        & Desempe\~{n}o de detecci\'{o}n y protocolo con que se mide (m\'{e}tricas, procedencia de etiquetas, umbral, unidad de evaluaci\'{o}n, partici\'{o}n) \\
Contexto         & Sistemas inteligentes de transporte; estudios revisados por pares publicados entre 2020 y 2026 \\
\bottomrule
\end{tabularx}
\end{table}

Para evitar que el alcance declarado difiera del corpus, el t\'{i}tulo, las preguntas, el PICOC, la cadena de b\'{u}squeda y los criterios se alinearon expl\'{i}citamente (Tabla~\ref{tab:alineamiento}).

\begin{table}[H]
\centering
\small
\caption{Alineamiento del alcance entre los componentes del protocolo.}
\label{tab:alineamiento}
\begin{tabularx}{\linewidth}{lYYY}
\toprule
\textbf{Componente} & \textbf{Entorno vial} & \textbf{Tipo de dato} & \textbf{M\'{e}todo} \\
\midrule
T\'{i}tulo y RQ1 & Urbano y autopista & Sensores, trayectorias y video & DL \\
PICOC (P, I) & Redes urbanas y autopistas & Sensores, GPS, movilidad agregada, video & DL como m\'{e}todo principal \\
Cadena (bloque C) & \textit{road, urban, vehicular, highway} o \textit{freeway traffic} & \textit{traffic flow, speed, data, video, surveillance} o \textit{camera}; \textit{vehicle trajector*}; \textit{taxi*} & Bloque B (arquitecturas de DL) \\
Criterios & CI-4, CE-2 & CI-4 & CI-5, CE-4 \\
\bottomrule
\end{tabularx}
\end{table}

\subsection{Fuentes de informaci\'{o}n y estrategia de b\'{u}squeda}
\label{sec:busqueda}

Se consultaron cinco bases bibliogr\'{a}ficas: Scopus, Web of Science Core Collection, IEEE Xplore, ScienceDirect y ACM Digital Library. Se incorpor\'{o} IEEE Xplore por su peso en el \'{a}rea de transporte inteligente y visi\'{o}n por computador. MDPI y Taylor \& Francis Online, que figuraban en la versi\'{o}n anterior del protocolo, se retiraron como fuentes porque son plataformas editoriales y no bases bibliogr\'{a}ficas. Sus revistas indexadas (p.~ej., \textit{Sensors} o \textit{Applied Sciences}) se recuperan a trav\'{e}s de Scopus y Web of Science, de modo que su retiro no excluye esos art\'{i}culos. Se us\'{o} Parsif.al para gestionar los registros.

La cadena se construy\'{o} a partir del PICOC con tres bloques unidos por AND: (A) t\'{e}rminos de detecci\'{o}n de anomal\'{i}as, incidentes o accidentes; (B) t\'{e}rminos de DL y de sus arquitecturas; y (C) t\'{e}rminos de tr\'{a}fico vehicular, incluidos los de video y vigilancia. La cadena anterior exig\'{i}a adem\'{a}s \textit{time series} o \textit{spatio-temporal}, lo que recuperaba solo parcialmente los estudios basados en video. Ese bloque se elimin\'{o} al ampliar el alcance. En todas las bases se busc\'{o} en t\'{i}tulo, resumen y palabras clave, salvo en ACM (resumen) e IEEE Xplore (\textit{All Metadata}). Se filtr\'{o} por a\~{n}o de publicaci\'{o}n (2020--2026), por tipo de documento (art\'{i}culo de revista o de conferencia) y por idioma (ingl\'{e}s o espa\~{n}ol). En Web of Science y ScienceDirect estos filtros se aplicaron en la interfaz. Scopus, IEEE Xplore y ACM se exportaron sin filtros y los mismos filtros se aplicaron con un script sobre los metadatos exportados (\NFueraFiltros{} registros fuera de periodo o de tipo de documento), con un registro de cada exclusi\'{o}n. El acceso abierto no se filtr\'{o} en las interfaces, porque cada base lo marca de forma distinta y ScienceDirect y ACM no lo exportan; se verific\'{o} registro por registro como parte del criterio CI-2 (Secci\'{o}n~\ref{sec:metodo}). Las cadenas exactas por base, con sus campos y filtros, se presentan en el Anexo~\ref{ap:protocolo}; la fecha de ejecuci\'{o}n y el n\'{u}mero de registros de cada base se reportan en la Tabla~\ref{tab:bases}. Como el periodo incluye el a\~{n}o en curso, los resultados de 2026 son parciales hasta la fecha de b\'{u}squeda.

La cadena fue elaborada por el autor; no fue validada por un especialista en informaci\'{o}n externo, lo que se discute en la Secci\'{o}n~\ref{sec:amenazas}. La cadena se valid\'{o} con un conjunto de control formado por los estudios incluidos: se comprob\'{o} de forma autom\'{a}tica si el t\'{i}tulo, el resumen y las palabras clave de cada uno satisfacen los tres bloques (script \texttt{prueba\_sensibilidad\_cadena.py}). La primera formulaci\'{o}n de la versi\'{o}n 2 recuperaba 13 de 21 estudios. Se a\~{n}adieron sin\'{o}nimos (\textit{anomal*}, \textit{abnormal*}, \textit{road closure}, \textit{expressway}, \textit{vehicle}, entre otros) y la cadena recuper\'{o} 19 de 21. Una prueba en Scopus mostr\'{o} que los t\'{e}rminos sueltos del bloque C (\textit{vehicle*}, \textit{highway*}) tra\'{i}an mucho ruido (veh\'{i}culos a\'{e}reos no tripulados, veh\'{i}culos el\'{e}ctricos). Por eso, en la versi\'{o}n definitiva (v2.1) se sustituyeron por frases ancladas al tr\'{a}fico vial (\textit{traffic speed}, \textit{traffic monitoring}, \textit{vehicle trajector*}, entre otras), y la cadena mantiene la recuperaci\'{o}n de \NsensibOK{} de \Nestudios{}. Los dos restantes (\IDsensibNo) no mencionan ning\'{u}n t\'{e}rmino de DL en el t\'{i}tulo, el resumen ni las palabras clave. Se identificaron en la b\'{u}squeda de la versi\'{o}n 1 en ACM Digital Library, que no restring\'{i}a el campo y por tanto buscaba en el texto completo, y se reportan en la rama ``otros m\'{e}todos'' del diagrama PRISMA (Figura~\ref{fig:prisma}).

\begin{table}[H]
\centering
\small
\caption{Bases consultadas: fecha de ejecuci\'{o}n, campos y registros recuperados.}
\label{tab:bases}
\begin{tabularx}{\linewidth}{lYcc}
\toprule
\textbf{Base} & \textbf{Campos y filtros} & \textbf{Fecha} & \textbf{Registros} \\
\midrule
Scopus & TITLE-ABS-KEY; 2020--2026; art\'{i}culo y conferencia; ingl\'{e}s/espa\~{n}ol & \FScopus & \NScopus \\
Web of Science & TS (t\'{i}tulo, resumen, palabras clave y Keywords Plus); mismos filtros & \FWoS & \NWoS \\
IEEE Xplore & \textit{All Metadata}; tres consultas (l\'{i}mite de 25 t\'{e}rminos); revistas y conferencias & \FIEEE & \NIEEE \\
ScienceDirect & T\'{i}tulo, resumen y palabras clave; cinco consultas (l\'{i}mite de 8 operadores, sin comodines) & \FScienceDirect & \NScienceDirect \\
ACM Digital Library & Resumen; art\'{i}culos de investigaci\'{o}n & \FACM & \NACM \\
\midrule
\textbf{Total} & & & \Pidentificados \\
\bottomrule
\end{tabularx}
\end{table}

\subsection{Criterios de elegibilidad}

\begin{table}[H]
\centering
\small
\caption{Criterios de inclusi\'{o}n (CI) y exclusi\'{o}n (CE), con su fundamento en el protocolo.}
\label{tab:criterios}
\begin{tabularx}{\linewidth}{lYP{3.4cm}}
\toprule
\textbf{C\'{o}d.} & \textbf{Criterio} & \textbf{Fundamento} \\
\midrule
\multicolumn{3}{l}{\cellcolor{azulClaro!40}\textbf{Inclusi\'{o}n}} \\
CI-1 & Art\'{i}culo de revista o de conferencia revisado por pares & Calidad m\'{i}nima de la evidencia \\
CI-2 & En ingl\'{e}s o espa\~{n}ol, con texto completo en acceso abierto & Viabilidad y reproducibilidad: la extracci\'{o}n exige el texto completo y cualquier lector debe poder verificarla sin suscripci\'{o}n \\
CI-3 & Publicado entre 2020 y 2026 & Contexto (C): periodo del DL reciente; vac\'{i}o de las revisiones previas \\
CI-4 & Aborda la detecci\'{o}n de anomal\'{i}as en tr\'{a}fico vehicular (urbano o de autopista) a partir de sensores, trayectorias, datos de movilidad o video & Poblaci\'{o}n (P) y alcance; RQ3, RQ4 \\
CI-5 & Usa \textit{Deep Learning} como m\'{e}todo principal de la detecci\'{o}n; en modelos h\'{i}bridos, la componente profunda aprende la representaci\'{o}n sobre la que se decide la anomal\'{i}a & Intervenci\'{o}n (I); RQ1 \\
CI-6 & Reporta una evaluaci\'{o}n cuantitativa de la tarea de detecci\'{o}n (no solo de una tarea auxiliar como pron\'{o}stico o reconocimiento de acciones) & Resultado (O); RQ2 y RQ6 \\
\midrule
\multicolumn{3}{l}{\cellcolor{rojoAlerta!10}\textbf{Exclusi\'{o}n}} \\
CE-1 & Duplicado entre bases, o versi\'{o}n menos completa del mismo estudio (p.~ej., conferencia ampliada luego en revista); se conserva la m\'{a}s completa & Evitar doble conteo \\
CE-2 & No trata de tr\'{a}fico vehicular por carretera (p.~ej., tr\'{a}fico de red, ciberseguridad del veh\'{i}culo, peatones, transporte mar\'{i}timo, a\'{e}reo o ferroviario, estado del pavimento, monitoreo del conductor) & Poblaci\'{o}n (P) \\
CE-3 & Solo pronostica o clasifica tr\'{a}fico, o estima riesgo a largo plazo, sin detectar ni anticipar anomal\'{i}as & Resultado (O); objetivo de la RSL \\
CE-4 & La detecci\'{o}n no se basa en DL (SVM, ARIMA, reglas o umbrales sin componente profunda entrenada) & Intervenci\'{o}n (I) \\
CE-5 & Revisi\'{o}n, \textit{survey}, cap\'{i}tulo, tesis, resumen, p\'{o}ster, editorial o \textit{preprint}; un \textit{preprint} se incluye solo si existe su versi\'{o}n revisada por pares, y entonces se cita esta \'{u}ltima & Solo estudios primarios revisados por pares \\
CE-6 & No describe el modelo, los datos o los resultados de forma suficiente para la extracci\'{o}n & Viabilidad de la extracci\'{o}n (Tabla~\ref{tab:formulario}) \\
\bottomrule
\end{tabularx}
\end{table}

Los criterios se aplicaron como un \'{a}rbol de decisi\'{o}n en orden fijo (CE-1, CE-5, CI-3, CI-2, CE-2, CE-3, CE-4 y, a texto completo, CI-6 y CE-6), de modo que cada exclusi\'{o}n se asigna al primer criterio que no se cumple. Las definiciones operativas, con ejemplos de inclusi\'{o}n y de exclusi\'{o}n, se recogen en el manual de cribado del material suplementario. Los casos l\'{i}mite (estudios en los que el DL no decide directamente la anomal\'{i}a o en los que la evaluaci\'{o}n cuantitativa corresponde a otra tarea) se resolvieron con CI-5 y CI-6, y la decisi\'{o}n se documenta estudio por estudio en el Anexo~\ref{ap:limite}.

\subsection{Proceso de selecci\'{o}n}

Los registros exportados de las cinco bases se unificaron y se deduplicaron con un script (coincidencia de DOI o de t\'{i}tulo normalizado con a\~{n}o igual o contiguo), cuyo registro de duplicados forma parte del material suplementario; los registros \'{u}nicos se importaron en Parsif.al. Al ser una revisi\'{o}n de autor \'{u}nico, la selecci\'{o}n la realiz\'{o} un solo revisor (\Revisor), con tres medidas para controlar el sesgo de selecci\'{o}n. Primero, un asistente de cribado basado en reglas transparentes (derivadas de la cadena de b\'{u}squeda y del \'{a}rbol de criterios) pre-clasific\'{o} cada registro como ``excluir por regla'', ``probable inclusi\'{o}n'' o ``revisar''. Solo propuso exclusiones por criterios objetivos (tipo de documento, a\~{n}o, idioma o dominio claramente no vial), y el revisor verific\'{o} cada una (\NExclReglaConfirmadas{} de \NExclRegla{} confirmadas); Adem\'{a}s, un asistente de IA generativa (Claude, Anthropic) ley\'{o} el t\'{i}tulo y el resumen de cada registro, aplic\'{o} el mismo \'{a}rbol de criterios y propuso una decisi\'{o}n con su criterio y una justificaci\'{o}n breve. El revisor verific\'{o} cada propuesta leyendo el t\'{i}tulo y el resumen y tom\'{o} la decisi\'{o}n final; cambi\'{o} \NCambiosIATA{} propuestas (acuerdo revisor--IA de \AcuerdoIATA{}, $\kappa$~=~\KappaIATA). Las propuestas originales de la IA se conservan en el material suplementario. El acuerdo entre el revisor y el asistente por reglas fue de \AcuerdoAsistenteTA{} ($\kappa$~=~\KappaAsistenteTA). El requisito de acceso abierto de CI-2 se verific\'{o} para cada registro el \FechaOA{} con OpenAlex, que integra los datos de Unpaywall, y con el indicador de acceso abierto de Scopus, Web of Science e IEEE Xplore. Se consider\'{o} de acceso abierto todo texto completo legible sin suscripci\'{o}n (\textit{gold}, \textit{diamond}, \textit{hybrid}, \textit{bronze} o \textit{green}) seg\'{u}n cualquiera de las fuentes; \NexclOA{} registros se excluyeron por no cumplirlo. Segundo, en t\'{i}tulo y resumen se excluy\'{o} solo cuando el incumplimiento era claro, y los casos dudosos pasaron a texto completo. A texto completo, un script extrajo de cada PDF las frases candidatas (m\'{e}tricas de detecci\'{o}n, m\'{e}todo propuesto y t\'{e}rminos de DL) con su p\'{a}gina, y el revisor decidi\'{o} leyendo el art\'{i}culo; cada exclusi\'{o}n se registr\'{o} con su criterio y con la cita que la justifica. Tercero, la consistencia del revisor se midi\'{o} mediante un re-cribado ciego, al menos siete d\'{i}as despu\'{e}s, de una muestra aleatoria: $\kappa$ intra-revisor~=~\KappaRecribadoTA{} ($n=\NRecribadoTA$) en t\'{i}tulo y resumen y \KappaRecribadoTC{} ($n=\NRecribadoTC$) en texto completo, con el coeficiente kappa de Cohen~\cite{cohen1960coefficient}. Las decisiones finales se consignaron en Parsif.al y los motivos de exclusi\'{o}n a texto completo se reportan en la Figura~\ref{fig:prisma}.

\subsection{Extracci\'{o}n de datos}

Los datos de cada estudio se extrajeron con el formulario de la Tabla~\ref{tab:formulario}. La extracci\'{o}n inicial se hizo con asistencia de inteligencia artificial (lectura del texto completo de cada estudio) y el autor la verific\'{o} contra el PDF, corrigiendo las discrepancias; las fichas individuales conservan la evidencia de cada dato. Adem\'{a}s de los campos descriptivos, se extrajeron la procedencia de las etiquetas, la proporci\'{o}n de anomal\'{i}as, la unidad de evaluaci\'{o}n, el criterio de umbral, la partici\'{o}n de los datos y las l\'{i}neas base, porque son los elementos que determinan la comparabilidad entre estudios. La matriz completa (un registro por estudio) se publica como material suplementario y todas las cifras de este art\'{i}culo se generan a partir de ella con un script (\texttt{generar\_conteos.py}).

\begin{table}[H]
\centering
\small
\caption{Formulario de extracci\'{o}n de datos.}
\label{tab:formulario}
\begin{tabularx}{\linewidth}{lY}
\toprule
\textbf{Campo} & \textbf{Contenido} \\
\midrule
Identificaci\'{o}n & ID, autores, a\~{n}o, t\'{i}tulo, revista o conferencia, DOI, base de recuperaci\'{o}n \\
Contexto & Entorno (urbano, autopista), pa\'{i}s o regi\'{o}n de los datos \\
Datos & Tipo de dato, variables de entrada, datasets (real, simulado o mixto; p\'{u}blico o no) \\
Etiquetas & Procedencia (registros reales, anotaci\'{o}n experta, benchmark de terceros, inyectadas), proporci\'{o}n de anomal\'{i}as \\
Modelo & Familias de arquitectura, h\'{i}brido o no, multimodal o no, paradigma (supervisado, semisupervisado, no supervisado) y mecanismo (predicci\'{o}n, reconstrucci\'{o}n, clasificaci\'{o}n, similitud) \\
Protocolo & Partici\'{o}n entrenamiento/prueba, unidad de evaluaci\'{o}n, criterio de umbral \\
Resultados & M\'{e}tricas, resultado principal, l\'{i}neas base y ventaja frente a la mejor \\
Anomal\'{i}a & Categor\'{i}a principal y secundaria \\
Limitaciones & Limitaciones declaradas por los autores (codificadas en seis categor\'{i}as) \\
Reproducibilidad & Disponibilidad de c\'{o}digo y datos \\
\bottomrule
\end{tabularx}
\end{table}

\subsection{Evaluaci\'{o}n de la calidad metodol\'{o}gica}

La versi\'{o}n anterior del instrumento repet\'{i}a los criterios de elegibilidad, de modo que todo estudio incluido los cumpl\'{i}a por definici\'{o}n. Se sustituy\'{o} por un instrumento de ocho \'{i}tems, adaptado de la lista de calidad de Kitchenham y Charters~\cite{Kitchenham2007} y de las recomendaciones REFORMS para estudios basados en aprendizaje autom\'{a}tico~\cite{kapoor2024reforms}. Los \'{i}tems eval\'{u}an riesgos de sesgo propios de este tipo de estudios (Tabla~\ref{tab:cqi}). Cada \'{i}tem se punt\'{u}a como S\'{i} (1), Parcial (0,5) o No (0), con un m\'{a}ximo de 8. La puntuaci\'{o}n \textbf{no} se us\'{o} para excluir estudios: se reporta por \'{i}tem y por estudio, y se usa en la s\'{i}ntesis para contrastar los hallazgos principales con el subconjunto de calidad alta ($\geq$\,\QAumbralAlta). Para cada estudio se propuso un puntaje por \'{i}tem con la evidencia del texto (asistido), y el autor lo verific\'{o} leyendo el art\'{i}culo, confirmando o corrigiendo cada puntaje con su justificaci\'{o}n; el acuerdo entre el puntaje propuesto y el verificado fue de \AcuerdoCalidad.

\begin{table}[H]
\centering
\small
\caption{Instrumento de evaluaci\'{o}n de calidad metodol\'{o}gica (QA).}
\label{tab:cqi}
\begin{tabularx}{\linewidth}{lP{3.6cm}Y}
\toprule
\textbf{\'{I}tem} & \textbf{Criterio} & \textbf{S\'{i} (1) / Parcial (0,5) / No (0)} \\
\midrule
QA1 & Procedencia de etiquetas & 1: registros reales o anotaci\'{o}n verificada; 0,5: mixtas, autom\'{a}ticas o verificaci\'{o}n parcial; 0: solo inyectadas o no reportado \\
QA2 & Partici\'{o}n sin fuga & 1: partici\'{o}n temporal o prueba independiente descrita; 0,5: partici\'{o}n aleatoria o descripci\'{o}n incompleta; 0: no reportada \\
QA3 & L\'{i}neas base comparables & 1: dos o m\'{a}s l\'{i}neas base con el mismo protocolo, incluida una del estado del arte o no DL; 0,5: una sola o solo ablaci\'{o}n; 0: ninguna \\
QA4 & Variabilidad & 1: varias ejecuciones con dispersi\'{o}n, intervalos o pruebas de significancia; 0,5: parcial; 0: una sola ejecuci\'{o}n \\
QA5 & C\'{o}digo y datos & 1: ambos p\'{u}blicos; 0,5: uno de los dos; 0: ninguno \\
QA6 & M\'{a}s de un escenario & 1: dos o m\'{a}s datasets, ciudades o v\'{i}as independientes; 0,5: varios subconjuntos de un dataset; 0: un \'{u}nico escenario \\
QA7 & Protocolo expl\'{i}cito & 1: reporta la unidad de evaluaci\'{o}n y el criterio de umbral o decisi\'{o}n; 0,5: solo uno; 0: ninguno \\
QA8 & Limitaciones & 1: secci\'{o}n o p\'{a}rrafo expl\'{i}cito; 0,5: menci\'{o}n breve; 0: ninguna \\
\bottomrule
\end{tabularx}
\end{table}

\subsection{S\'{i}ntesis}

Los estudios se agruparon por familia de arquitectura, tipo de dato, categor\'{i}a de anomal\'{i}a y periodo de publicaci\'{o}n. Para sustentar las tendencias se construyeron tablas cruzadas (tipo de dato $\times$ familia, anomal\'{i}a $\times$ familia y periodo $\times$ familia). Se identificaron los conjuntos de datos usados por m\'{a}s de un estudio para comparar su desempe\~{n}o bajo el protocolo de cada uno. No se realiz\'{o} metaan\'{a}lisis porque los estudios difieren en tarea, datos, etiquetas, unidad de evaluaci\'{o}n y m\'{e}trica, y ning\'{u}n par de estudios comparte a la vez dataset y protocolo. Por la misma raz\'{o}n no se aplic\'{o} GRADE para valorar la certeza de la evidencia ni se evalu\'{o} formalmente el sesgo de publicaci\'{o}n (PRISMA, \'{i}tems 14 y 15). Como aproximaci\'{o}n, se registr\'{o} si cada estudio compara contra l\'{i}neas base y si reporta resultados desfavorables para su propuesta (Secci\'{o}n~\ref{sec:rq6}).

%=============================================================
\section{Resultados}
\label{sec:resultados}
%=============================================================

\subsection{Selecci\'{o}n de estudios}

La Figura~\ref{fig:prisma} resume el proceso de selecci\'{o}n seg\'{u}n PRISMA 2020. De \Pidentificados{} registros identificados en las cinco bases se eliminaron \Pduplicados{} duplicados y se cribaron \Pcribados{} por t\'{i}tulo y resumen. Se evaluaron \Pevaluados{} textos completos, de los que se excluyeron \PexclTC{} con los motivos indicados en la figura. A los estudios procedentes de las bases se sumaron \POtrosIncluidos{} estudios identificados en la b\'{u}squeda de la versi\'{o}n 1 en ACM Digital Library (\IDsensibNo). En total se incluyeron \Nestudios{} estudios. \NexcluidosLimite{} estudios incluidos en la versi\'{o}n anterior de esta revisi\'{o}n o en la planilla de Parsif.al se excluyeron al aplicar los criterios actualizados; estas exclusiones y las \NincluidosLimite{} inclusiones dudosas se justifican estudio por estudio en el Anexo~\ref{ap:limite}.

\begin{figure}[H]
\centering
\small
\begin{tikzpicture}[
  caja/.style={draw, rectangle, rounded corners=2pt, align=left, text width=5.0cm, minimum height=1.0cm, font=\footnotesize, fill=white},
  excl/.style={draw, rectangle, rounded corners=2pt, align=left, text width=4.4cm, font=\footnotesize, fill=grisClaro},
  otro/.style={draw, rectangle, rounded corners=2pt, align=left, text width=2.9cm, font=\footnotesize, fill=azulClaro!40},
  fase/.style={draw, fill=azulPrincipal, text=white, font=\footnotesize\bfseries, minimum width=0.55cm, rotate=90, align=center},
  flecha/.style={-{Stealth[length=2mm]}, thick}]
\node[caja] (id) {Registros identificados en bases ($n=\Pidentificados$):\\Scopus ($n=\NScopus$), WoS ($n=\NWoS$),\\IEEE Xplore ($n=\NIEEE$), ScienceDirect ($n=\NScienceDirect$), ACM DL ($n=\NACM$)};
\node[excl, right=0.5cm of id] (dup) {Eliminados antes del cribado:\\duplicados ($n=\Pduplicados$)};
\node[caja, below=0.7cm of id] (crib) {Registros cribados por t\'{i}tulo y resumen\\($n=\Pcribados$)};
\node[excl, right=0.5cm of crib] (exta) {Registros excluidos ($n=\PexclTA$)};
\node[caja, below=0.7cm of crib] (bus) {Informes buscados a texto completo\\($n=\Pbuscados$)};
\node[excl, right=0.5cm of bus] (norec) {Informes no recuperados ($n=\PnoRecuperados$)};
\node[caja, below=0.7cm of bus] (eleg) {Informes evaluados para elegibilidad\\($n=\Pevaluados$)};
\node[excl, right=0.5cm of eleg] (extc) {Informes excluidos ($n=\PexclTC$):\\CE-2 no es tr\'{a}fico vehicular ($n=\PexclCEdos$)\\CE-3 solo pron\'{o}stico ($n=\PexclCEtres$)\\CE-4 sin DL principal ($n=\PexclCEcuatro$)\\CE-5 tipo de documento ($n=\PexclCEcinco$)\\CI-6 sin evaluaci\'{o}n de la detecci\'{o}n ($n=\PexclCIseis$)\\CE-6 informaci\'{o}n insuficiente ($n=\PexclCEseis$)\\Otros: CE-1, CI-2, CI-3 ($n=\PexclOtros$)};
\node[caja, below=1.6cm of eleg] (inc) {\textbf{Estudios incluidos en la revisi\'{o}n ($n=\Nestudios$)}};
\node[otro, left=0.4cm of eleg] (otros) {Otros m\'{e}todos: b\'{u}squeda v1 en ACM DL a texto completo ($n=\POtrosIdentificados$; incluidos $n=\POtrosIncluidos$)};
\draw[flecha] (id) -- (crib);
\draw[flecha] (id) -- (dup);
\draw[flecha] (crib) -- (bus);
\draw[flecha] (crib) -- (exta);
\draw[flecha] (bus) -- (eleg);
\draw[flecha] (bus) -- (norec);
\draw[flecha] (eleg) -- (inc);
\draw[flecha] (eleg) -- (extc);
\draw[flecha] (otros) |- (inc);
\node[fase, left=4.0cm of id.west, anchor=south] {Identificaci\'{o}n};
\node[fase, left=4.0cm of bus.west, anchor=south] {Cribado};
\node[fase, left=4.0cm of inc.west, anchor=south] {Incluidos};
\end{tikzpicture}
\caption{Diagrama de flujo PRISMA 2020 de la selecci\'{o}n de estudios.}
\label{fig:prisma}
\end{figure}

\subsection{Caracter\'{i}sticas de los estudios}

Los \Nestudios{} estudios se publicaron entre \AnioMin{} y \AnioMax{}: \AnioA{} en 2020, \AnioB{} en 2021, \AnioC{} en 2022, \AnioD{} en 2023, \AnioE{} en 2024, \AnioF{} en 2025 y \AnioG{} en 2026 (a\~{n}o parcial, hasta la fecha de b\'{u}squeda). \NrecientesVeintidosVeintiseis{} estudios son de 2024 en adelante. Por entorno, \Nurbano{} estudios analizan tr\'{a}fico urbano, \Nautopista{} tr\'{a}fico de autopista, \Nambos{} ambos y \NentornoNE{} no lo especifican. Por tipo de dato, \NdatoVideo{} usan video o im\'{a}genes, \NdatoSensores{} series de sensores, \NdatoTrayectorias{} trayectorias, \NdatoMovilidad{} datos de movilidad agregada u origen-destino, \NdatoMultimodal{} combinan varias fuentes y \NdatoSimulado{} usan exclusivamente datos simulados. La Tabla~\ref{tab:caracteristicas} presenta las caracter\'{i}sticas y el resultado principal de cada estudio.

\begin{landscape}
\footnotesize
\setlength{\tabcolsep}{3pt}
\begin{longtable}{llcP{2.5cm}P{1.3cm}P{3.0cm}P{2.0cm}P{2.4cm}P{3.5cm}P{3.0cm}}
\caption{Caracter\'{i}sticas de los estudios incluidos.\label{tab:caracteristicas}}\\
\toprule
\textbf{ID} & \textbf{Ref.} & \textbf{A\~{n}o} & \textbf{Tipo de dato} & \textbf{Entorno} & \textbf{Datasets} & \textbf{Etiquetas} & \textbf{Anomal\'{i}a} & \textbf{Modelo} & \textbf{Resultado principal} \\
\midrule\endfirsthead
\toprule
\textbf{ID} & \textbf{Ref.} & \textbf{A\~{n}o} & \textbf{Tipo de dato} & \textbf{Entorno} & \textbf{Datasets} & \textbf{Etiquetas} & \textbf{Anomal\'{i}a} & \textbf{Modelo} & \textbf{Resultado principal} \\
\midrule\endhead
\bottomrule\endfoot
\tabinput{generado/caracteristicas.tex}
\end{longtable}
\end{landscape}

\subsection{Calidad metodol\'{o}gica}

La puntuaci\'{o}n media fue \QAmedia{} sobre 8 (desviaci\'{o}n est\'{a}ndar \QAde; rango \QAmin--\QAmax). El instrumento discrimina entre estudios, a diferencia de la versi\'{o}n anterior (Tabla~\ref{tab:calidad}). Los \'{i}tems con menor cumplimiento son la variabilidad entre ejecuciones (QA4: \QACuatrosi{} estudios con S\'{i}) y la disponibilidad de c\'{o}digo y datos (QA5: \QACincosi{}), seguidos de la procedencia verificada de las etiquetas (QA1: \QAUnosi{}). Los de mayor cumplimiento son las l\'{i}neas base (QA3: \QATressi{}) y el protocolo expl\'{i}cito (QA7: \QASietesi{}). \NaltaCalidad{} estudios alcanzan la calidad alta ($\geq$\,\QAumbralAlta): \IDaltaCalidad.

\begin{table}[H]
\centering
\small
\caption{Evaluaci\'{o}n de calidad por \'{i}tem y estudio (1 = S\'{i}, 0.5 = Parcial, 0 = No).}
\label{tab:calidad}
\begin{tabular}{lccccccccc}
\toprule
\textbf{ID} & \textbf{QA1} & \textbf{QA2} & \textbf{QA3} & \textbf{QA4} & \textbf{QA5} & \textbf{QA6} & \textbf{QA7} & \textbf{QA8} & \textbf{Total} \\
\midrule
\tabinput{generado/calidad.tex}
\bottomrule
\end{tabular}
\end{table}

\subsection{RQ1: T\'{e}cnicas de \textit{Deep Learning} y su relaci\'{o}n con el tipo de dato}

La Tabla~\ref{tab:familias} muestra la presencia de cada familia de arquitectura. Predominan las redes recurrentes (\NfamRec{} estudios), las convolucionales (\NfamConv{}) y los mecanismos de atenci\'{o}n o Transformer (\NfamTransf{}), seguidos de las redes de grafos (\NfamGrafo{}), los autoencoders (\NfamAE{}) y los modelos generativos (\NfamGen{}). La mayor\'{i}a de los estudios (\Nhibrido{} de \Nestudios) combina dos o m\'{a}s familias, por lo que la clasificaci\'{o}n en una sola familia, usada en la versi\'{o}n anterior, ocultaba esta composici\'{o}n. Se distingue entre modelos \textit{h\'{i}bridos}, que combinan arquitecturas, y \textit{multimodales}, que integran varias modalidades de datos (\Nmultimodal{} estudios). Seg\'{u}n el mecanismo de detecci\'{o}n, \NenfoqueClasif{} estudios clasifican de forma supervisada, \NenfoquePred{} detectan por error de predicci\'{o}n y \NenfoqueRecon{} por error de reconstrucci\'{o}n.

% Generado por suplementos/generar_conteos.py. No editar a mano.
\begingroup
\small
\begin{xltabular}{\linewidth}{lccX}
\caption{Familias de arquitectura presentes en los estudios (RQ1).}\label{tab:familias}\\
\toprule
\textbf{Familia} & \textbf{n} & \textbf{\%} & \textbf{Estudios} \\
\midrule\endfirsthead
\toprule
\textbf{Familia} & \textbf{n} & \textbf{\%} & \textbf{Estudios} \\
\midrule\endhead
\bottomrule\endfoot
Recurrentes (LSTM, GRU, Bi-LSTM) & 34 & 47~\% & S02, S03, S04, S05, S09, S10, S13, S14, S15, S16, S18, S20, S21, S22, S23, S24, S27, S28, S29, S30, S32, S33, S34, S38, S40, S41, S46, S47, S50, S62, S65, S66, S69, S72 \\
Convolucionales (CNN, TCN, 3D-CNN) & 43 & 60~\% & S01, S02, S03, S04, S08, S09, S17, S19, S22, S25, S26, S27, S28, S29, S30, S32, S36, S37, S39, S40, S42, S43, S44, S45, S46, S47, S49, S50, S52, S53, S55, S56, S57, S58, S60, S61, S64, S66, S67, S68, S70, S71, S72 \\
Autoencoders y reconstrucci\'{o}n & 13 & 18~\% & S09, S11, S14, S15, S18, S20, S23, S31, S32, S47, S50, S60, S63 \\
Grafos (GNN, GCN, GAT) & 11 & 15~\% & S12, S18, S31, S32, S36, S40, S48, S51, S57, S58, S63 \\
Transformers y atenci\'{o}n & 34 & 47~\% & S07, S12, S14, S17, S18, S22, S26, S27, S29, S34, S35, S39, S40, S42, S44, S45, S46, S47, S48, S49, S50, S51, S53, S54, S57, S58, S59, S61, S62, S63, S68, S69, S70, S71 \\
Generativos (GAN, VAE, difusi\'{o}n) & 8 & 11~\% & S11, S14, S22, S23, S34, S47, S60, S63 \\
\midrule
Combinan dos o m\'{a}s familias (h\'{i}bridos) & 36 & 50~\% & S02, S03, S04, S09, S12, S14, S17, S18, S20, S22, S23, S24, S26, S27, S29, S30, S31, S32, S34, S35, S36, S39, S40, S42, S46, S47, S48, S50, S53, S57, S58, S61, S62, S63, S68, S69 \\
Integran varias modalidades de datos (multimodales) & 11 & 15~\% & S22, S28, S36, S39, S40, S50, S53, S58, S59, S62, S70 \\
\end{xltabular}
\vspace{-6pt}{\footnotesize\raggedright Un estudio puede combinar varias familias. ``H\'{i}brido'' se refiere a la arquitectura y ``multimodal'' a los datos de entrada.\par}
\endgroup


El cruce entre tipo de dato y familia (Tabla~\ref{tab:cruce_dato}) muestra que la arquitectura depende sobre todo del dato disponible. En video, \XVideoConv{} de \NdatoVideo{} estudios usan redes convolucionales, frecuentemente combinadas con redes recurrentes para la dimensi\'{o}n temporal (\XVideoRec{}), y ninguno usa grafos. Los \NgrafoNoVideo{} estudios con redes de grafos trabajan con sensores, movilidad agregada, trayectorias, datos multimodales o simulados, es decir, con datos que tienen una estructura de red vial expl\'{i}cita~\cite{hu2022detecting,ren2024graph,zhang2022automatic,park2025graph}. Los autoencoders aparecen sobre todo cuando no hay etiquetas suficientes: con matrices origen-destino~\cite{hu2022detecting,ren2024graph}, con trayectorias~\cite{wang2024multiscale} o con video entrenado solo con datos normales~\cite{pawar2022deep}.

% Generado por suplementos/generar_conteos.py. No editar a mano.
\begin{table}[H]
\centering
\small
\caption{Cruce tipo de dato $\times$ familia de arquitectura (n.\textsuperscript{o} de estudios).}\label{tab:cruce_dato}
\begin{tabularx}{\linewidth}{Xccccccc}
\toprule
\textbf{} & \textbf{Rec.} & \textbf{Conv.} & \textbf{AE} & \textbf{Grafo} & \textbf{Transf.} & \textbf{Gen.} & \textbf{Total} \\
\midrule
Series de sensores & 11 & 4 & 2 & 4 & 4 & 1 & 17 \\
Movilidad agregada/OD & 2 & 2 & 1 & 1 & 0 & 0 & 2 \\
Trayectorias & 9 & 2 & 6 & 3 & 6 & 4 & 13 \\
Video e im\'{a}genes & 7 & 26 & 3 & 0 & 15 & 2 & 29 \\
Multimodal & 5 & 9 & 1 & 3 & 9 & 1 & 11 \\
\bottomrule
\end{tabularx}
\par\vspace{2pt}{\footnotesize\raggedright Un estudio puede usar varias familias, por lo que las filas no suman el total.\par}
\end{table}


El cruce por periodo (Tabla~\ref{tab:cruce_anio}) sustenta la tendencia temporal con datos y no solo con impresiones. Ninguno de los \NperiodoAntes{} estudios de 2020--2021 usa atenci\'{o}n o Transformers, mientras que \XRecienteTransf{} de los \NperiodoReciente{} estudios de 2024--2026 los incorporan. Las redes de grafos pasan de \XAntesGrafo{} estudio en 2020--2021 a \XRecienteGrafo{} en 2024--2026. Las redes convolucionales y recurrentes siguen presentes en todos los periodos. Dado el tama\~{n}o del corpus, estas diferencias son descriptivas y no permiten inferir una tendencia estad\'{i}stica.

% Generado por suplementos/generar_conteos.py. No editar a mano.
\begin{table}[H]
\centering
\small
\caption{Cruce periodo de publicaci\'{o}n $\times$ familia de arquitectura.}\label{tab:cruce_anio}
\begin{tabularx}{\linewidth}{Xccccccc}
\toprule
\textbf{} & \textbf{Rec.} & \textbf{Conv.} & \textbf{AE} & \textbf{Grafo} & \textbf{Transf.} & \textbf{Gen.} & \textbf{Total} \\
\midrule
2020--2021 & 1 & 3 & 0 & 1 & 0 & 0 & 3 \\
2022--2023 & 4 & 3 & 3 & 2 & 2 & 1 & 7 \\
2024--2026 & 9 & 6 & 3 & 5 & 8 & 1 & 11 \\
\bottomrule
\end{tabularx}
\par\vspace{2pt}{\footnotesize\raggedright 2026 es parcial (b\'{u}squeda cerrada en la fecha indicada en la Secci\'{o}n~\ref{sec:busqueda}).\par}
\end{table}


\begin{hallazgo}
El tipo de dato condiciona la arquitectura: las redes convolucionales dominan en video y las de grafos solo aparecen cuando el dato tiene estructura de red vial. La mayor\'{i}a de los modelos son h\'{i}bridos, y desde 2024 la atenci\'{o}n o los Transformers est\'{a}n presentes en la mayor\'{i}a de los estudios.
\end{hallazgo}

\subsection{RQ2: M\'{e}tricas y protocolos de evaluaci\'{o}n}

\paragraph{M\'{e}tricas.} La Tabla~\ref{tab:metricas} muestra las m\'{e}tricas reportadas. La m\'{a}s frecuente es la \textit{accuracy} (\Nmetaccuracy{} estudios), seguida de recall (\Nmetrecall{}), AUC-ROC (\Nmetaucroc{}), precision (\Nmetprecision{}) y F1-score (\NmetFuno{}). Solo \Nmetprauc{} estudios reportan PR-AUC~\cite{chakraborty2023structural,wang2024multiscale}, pese a ser m\'{a}s informativa que la AUC-ROC con clases desbalanceadas. Las m\'{e}tricas propias de la detecci\'{o}n de incidentes (tasa de detecci\'{o}n, tasa de falsas alarmas y tiempo medio de detecci\'{o}n) aparecen en pocos estudios~\cite{fang2020deep,mehrannia2023deep,zhang2022automatic}. \Nmetoperativaslatenciafpsenergia{} estudios informan alguna m\'{e}trica operativa (latencia, FPS, tiempo de inferencia o energ\'{i}a). Varios estudios usan definiciones no est\'{a}ndar: por ejemplo, Mao et al.~\cite{mao2024detecting} llaman \textit{accuracy} a lo que corresponde a la precisi\'{o}n, y Park et al.~\cite{park2025graph} llaman \textit{recall} a una f\'{o}rmula de precisi\'{o}n.

% Generado por suplementos/generar_conteos.py. No editar a mano.
\begin{table}[H]
\centering
\small
\caption{M\'{e}tricas reportadas por los estudios (RQ2).}\label{tab:metricas}
\begin{tabularx}{\linewidth}{lccX}
\toprule
\textbf{M\'{e}trica} & \textbf{n} & \textbf{\%} & \textbf{Estudios} \\
\midrule
Accuracy & 43 & 60~\% & S02, S03, S04, S05, S10, S11, S12, S17, S19, S20, S21, S22, S25, S26, S27, S28, S29, S30, S33, S34, S36, S38, S40, S41, S43, S49, S50, S53, S54, S56, S58, S59, S60, S61, S62, S64, S65, S66, S67, S68, S70, S71, S72 \\
Precision & 52 & 72~\% & S02, S03, S05, S07, S08, S10, S11, S14, S15, S16, S19, S20, S22, S23, S24, S25, S26, S29, S31, S32, S33, S34, S35, S36, S38, S40, S41, S43, S44, S45, S47, S48, S49, S50, S51, S53, S54, S55, S56, S57, S59, S60, S61, S62, S63, S64, S65, S67, S68, S69, S71, S72 \\
Recall & 63 & 88~\% & S02, S03, S07, S08, S10, S11, S14, S15, S16, S17, S18, S19, S20, S21, S22, S23, S24, S25, S26, S27, S28, S29, S30, S31, S32, S33, S34, S35, S36, S37, S38, S39, S40, S41, S42, S43, S44, S45, S46, S47, S48, S49, S50, S51, S52, S54, S55, S56, S57, S58, S59, S60, S61, S62, S64, S65, S66, S67, S68, S69, S70, S71, S72 \\
F1-score & 60 & 83~\% & S02, S04, S05, S07, S09, S10, S11, S14, S15, S16, S17, S18, S19, S20, S22, S23, S24, S25, S26, S27, S28, S29, S30, S31, S32, S33, S35, S36, S37, S38, S39, S40, S41, S42, S43, S44, S47, S48, S49, S50, S51, S52, S53, S54, S55, S56, S58, S59, S60, S61, S62, S64, S65, S66, S67, S68, S69, S70, S71, S72 \\
AUC-ROC & 52 & 72~\% & S02, S05, S06, S07, S08, S09, S14, S15, S17, S18, S19, S20, S21, S22, S23, S24, S25, S27, S28, S29, S30, S31, S32, S35, S36, S37, S38, S39, S40, S41, S42, S43, S46, S47, S49, S50, S51, S53, S56, S57, S58, S59, S60, S61, S62, S63, S65, S66, S67, S69, S70, S72 \\
PR-AUC & 32 & 44~\% & S03, S09, S15, S18, S19, S20, S22, S23, S27, S28, S29, S30, S34, S37, S39, S42, S43, S45, S48, S49, S50, S51, S52, S54, S59, S61, S65, S66, S67, S69, S70, S72 \\
Tasa de detecci\'{o}n (DR) & 9 & 12~\% & S01, S05, S10, S12, S13, S15, S41, S62, S64 \\
FAR / FPR & 34 & 47~\% & S01, S05, S06, S10, S12, S13, S14, S15, S17, S19, S20, S21, S22, S27, S29, S35, S37, S39, S41, S43, S46, S49, S50, S51, S59, S61, S62, S63, S64, S65, S66, S67, S70, S72 \\
Tiempo medio de detecci\'{o}n & 12 & 17~\% & S01, S05, S13, S14, S19, S20, S34, S42, S43, S49, S54, S72 \\
MAE / RMSE / MAPE & 5 & 7~\% & S15, S31, S32, S63, S64 \\
IoU / mAP & 13 & 18~\% & S09, S17, S19, S27, S29, S37, S42, S43, S44, S45, S52, S55, S70 \\
Latencia / FPS / energ\'{i}a & 22 & 31~\% & S06, S10, S17, S19, S24, S31, S32, S34, S37, S46, S47, S50, S52, S56, S58, S59, S63, S66, S67, S68, S70, S71 \\
\bottomrule
\end{tabularx}
\par\vspace{2pt}{\footnotesize\raggedright Un estudio puede reportar varias m\'{e}tricas.\par}
\end{table}


\paragraph{Procedencia de las etiquetas.} En detecci\'{o}n de anomal\'{i}as, el origen de las etiquetas de prueba condiciona el resultado tanto como el modelo. Solo \NetiqReal{} estudios eval\'{u}an con registros reales de incidentes y \NetiqManual{} con anotaci\'{o}n manual propia (Tabla~\ref{tab:etiquetas}). \NetiqInyect{} estudios eval\'{u}an \'{u}nicamente con anomal\'{i}as inyectadas o sint\'{e}ticas~\cite{hu2022detecting,ren2024graph,wang2024multiscale,park2025graph}, \NetiqMixto{} con pseudo-etiquetas generadas autom\'{a}ticamente (por ejemplo, con Isolation Forest o con un autoencoder previo)~\cite{abdallah2025aipurban,duan2025multimodal,jadhav2026multistage} y \NetiqSin{} sin etiquetas, contrastando sus detecciones con el consenso de otros detectores~\cite{jia2025prediction}.

% Generado por suplementos/generar_conteos.py. No editar a mano.
\begingroup
\small
\begin{xltabular}{\linewidth}{lccX}
\caption{Origen de las etiquetas de anomal\'{i}a (RQ2, protocolo).}\label{tab:etiquetas}\\
\toprule
\textbf{Origen de las etiquetas} & \textbf{n} & \textbf{\%} & \textbf{Estudios} \\
\midrule\endfirsthead
\toprule
\textbf{Origen de las etiquetas} & \textbf{n} & \textbf{\%} & \textbf{Estudios} \\
\midrule\endhead
\bottomrule\endfoot
Registros reales de incidentes & 19 & 26~\% & S01, S05, S06, S11, S12, S15, S21, S35, S36, S38, S40, S41, S44, S53, S56, S58, S63, S65, S70 \\
Anotaci\'{o}n manual o experta & 9 & 12~\% & S02, S03, S07, S19, S24, S37, S43, S45, S61 \\
Benchmark anotado por terceros & 22 & 31~\% & S04, S09, S17, S22, S25, S26, S27, S29, S39, S42, S46, S47, S49, S52, S55, S59, S60, S64, S66, S67, S68, S71 \\
Anomal\'{i}as inyectadas o sint\'{e}ticas & 12 & 17~\% & S13, S14, S16, S20, S31, S32, S34, S48, S51, S57, S69, S72 \\
Mixto o autom\'{a}tico (pseudo-etiquetas) & 4 & 6~\% & S10, S28, S50, S62 \\
Sin etiquetas & 2 & 3~\% & S18, S30 \\
No reportado & 4 & 6~\% & S08, S23, S33, S54 \\
\end{xltabular}
\endgroup


\paragraph{Protocolo.} La Tabla~\ref{tab:protocolo} muestra qu\'{e} elementos del protocolo reporta cada estudio. Casi todos indican la unidad de evaluaci\'{o}n (\Nunidad{} de \Nestudios) y la partici\'{o}n de los datos (\Nparticion), aunque varios usan particiones aleatorias sobre series temporales, con riesgo de fuga de informaci\'{o}n (QA2). Del criterio de decisi\'{o}n, \NumbralExplicito{} estudios describen expl\'{i}citamente c\'{o}mo fijan el umbral y \NumbralNoAplica{} usan clasificadores que deciden sin umbral expl\'{i}cito o eval\'{u}an con m\'{e}tricas independientes del umbral. En varios casos el umbral se ajusta sobre los mismos eventos etiquetados que se usan para evaluar~\cite{cai2021road} o sobre una parte del conjunto de prueba~\cite{chakraborty2023structural}. Solo \Ntasa{} estudios reportan la proporci\'{o}n de anomal\'{i}as en los datos de evaluaci\'{o}n, un dato necesario para interpretar precision, F1 y PR-AUC.

% Generado por suplementos/generar_conteos.py. No editar a mano.
\begin{table}[H]
\centering
\small
\caption{Elementos del protocolo de evaluaci\'{o}n que reporta cada estudio (RQ2).}\label{tab:protocolo}
\begin{tabularx}{\linewidth}{lcX}
\toprule
\textbf{Elemento del protocolo reportado} & \textbf{n} & \textbf{Estudios} \\
\midrule
Tasa de anomal\'{i}as en los datos & 9 & S01, S02, S04, S05, S06, S11, S13, S14, S21 \\
Unidad de evaluaci\'{o}n (punto, ventana, frame, evento\ldots) & 20 & S01, S02, S03, S04, S05, S06, S07, S08, S09, S10, S11, S12, S13, S14, S15, S16, S18, S19, S20, S21 \\
Criterio de umbral o decisi\'{o}n & 19 & S01, S02, S03, S04, S05, S06, S07, S08, S09, S10, S11, S12, S15, S16, S17, S18, S19, S20, S21 \\
Partici\'{o}n entrenamiento/prueba & 18 & S01, S02, S03, S04, S05, S06, S08, S09, S10, S11, S13, S14, S16, S17, S18, S19, S20, S21 \\
\bottomrule
\end{tabularx}
\end{table}


\paragraph{Desempe\~{n}o sobre datasets compartidos.} Ning\'{u}n par de estudios comparte a la vez dataset, etiquetas y protocolo, lo que impide un metaan\'{a}lisis. Dos casos ilustran el problema. Primero, cuatro estudios usan datos de PeMS (PEMS-BAY, METR-LA, PEMS04/08 o estaciones de PeMS), pero cada uno define las anomal\'{i}as de forma distinta: pseudo-etiquetas de un autoencoder (F1 = 0,93 en PEMS-BAY)~\cite{abdallah2025aipurban}, registros de eventos combinados con patrones detectados autom\'{a}ticamente (F1 = 0,967)~\cite{duan2025multimodal}, ninguna etiqueta~\cite{jia2025prediction} y anomal\'{i}as generadas con Isolation Forest al 1~\% (F1 = 0,81)~\cite{jadhav2026multistage}. Las cifras no son comparables entre s\'{i}. Segundo, en NYC taxi con anomal\'{i}as inyectadas, Hu et al.~\cite{hu2022detecting} obtienen AUC entre 0,84 y 0,90 con Con-GAE, mientras que Ren et al.~\cite{ren2024graph} reportan AUC = 1,000 en todas las configuraciones, con un protocolo de inyecci\'{o}n distinto. Un AUC perfecto con anomal\'{i}as sint\'{e}ticas sugiere que la tarea de evaluaci\'{o}n es demasiado f\'{a}cil y no que el problema est\'{e} resuelto.

\begin{hallazgo}
La comparabilidad entre estudios est\'{a} limitada sobre todo por la procedencia de las etiquetas y por el protocolo, no por la arquitectura. Menos de un tercio de los estudios eval\'{u}a con etiquetas reales verificadas y pocos reportan la proporci\'{o}n de anomal\'{i}as o la variabilidad entre ejecuciones.
\end{hallazgo}

\subsection{RQ3: Tipos de anomal\'{i}as}

Seg\'{u}n la categor\'{i}a principal (Tabla~\ref{tab:categorias}), la m\'{a}s estudiada es accidentes e incidentes (\NcatA{} estudios, \PcatA), seguida de congesti\'{o}n o flujo an\'{o}malo (\NcatB, \PcatB), trayectorias o conducci\'{o}n peligrosa (\NcatC, \PcatC) y disrupciones y eventos extremos (\NcatD, \PcatD). Si se cuentan tambi\'{e}n las categor\'{i}as secundarias, los accidentes aparecen en \NcatAnyA{} estudios y la congesti\'{o}n en \NcatAnyB. A diferencia de la versi\'{o}n anterior, cada estudio se asigna a una sola categor\'{i}a principal, de modo que los porcentajes suman 100~\%.

% Generado por suplementos/generar_conteos.py. No editar a mano.
\begin{table}[H]
\centering
\small
\caption{Categor\'{i}as de anomal\'{i}a abordadas (RQ3).}\label{tab:categorias}
\begin{tabularx}{\linewidth}{lccXX}
\toprule
\textbf{Categor\'{i}a} & \textbf{n} & \textbf{\%} & \textbf{Estudios (principal)} & \textbf{Secundaria} \\
\midrule
Accidentes e incidentes & 9 & 43~\% & S01, S05, S06, S07, S10, S11, S17, S19, S20 & S02, S08, S21 \\
Congesti\'{o}n o flujo an\'{o}malo & 5 & 24~\% & S08, S12, S13, S18, S21 & S04, S16, S17, S19 \\
Trayectorias o conducci\'{o}n peligrosa & 4 & 19~\% & S02, S09, S14, S15 & S07 \\
Disrupciones y eventos extremos & 3 & 14~\% & S03, S04, S16 & S12, S18 \\
\midrule
Total & 21 & 100~\% &  &  \\
\bottomrule
\end{tabularx}
\end{table}


El cruce con la familia de arquitectura (Tabla~\ref{tab:cruce_anomalia}) refleja la relaci\'{o}n con el tipo de dato. Los accidentes se estudian sobre todo con video y redes convolucionales. La congesti\'{o}n y las disrupciones se estudian con series de sensores, movilidad agregada o trayectorias, y con modelos recurrentes o de grafos. Las trayectorias peligrosas concentran los modelos generativos y de reconstrucci\'{o}n~\cite{chakraborty2023structural,wang2024multiscale}.

% Generado por suplementos/generar_conteos.py. No editar a mano.
\begin{table}[H]
\centering
\small
\caption{Cruce categor\'{i}a principal de anomal\'{i}a $\times$ familia de arquitectura.}\label{tab:cruce_anomalia}
\begin{tabularx}{\linewidth}{Xccccccc}
\toprule
\textbf{} & \textbf{Rec.} & \textbf{Conv.} & \textbf{AE} & \textbf{Grafo} & \textbf{Transf.} & \textbf{Gen.} & \textbf{Total} \\
\midrule
Accidentes e incidentes & 16 & 29 & 5 & 4 & 19 & 3 & 43 \\
Congesti\'{o}n o flujo an\'{o}malo & 4 & 3 & 1 & 2 & 3 & 0 & 6 \\
Trayectorias o conducci\'{o}n peligrosa & 13 & 10 & 7 & 5 & 11 & 5 & 21 \\
Disrupciones y eventos extremos & 1 & 1 & 0 & 0 & 1 & 0 & 2 \\
\bottomrule
\end{tabularx}
\par\vspace{2pt}{\footnotesize\raggedright Clasificaci\'{o}n por categor\'{i}a principal (una por estudio).\par}
\end{table}


\subsection{RQ4: Conjuntos de datos}

La Tabla~\ref{tab:datos} agrupa los conjuntos de datos por tipo de dato. Los m\'{a}s reutilizados son los de PeMS (PEMS-BAY, METR-LA, PEMS04 y PEMS08) y los de taxis (NYC, Chicago, Porto y Chengdu). En video predominan los conjuntos espec\'{i}ficos de cada estudio o de competiciones (TNAD, VAID, IITH, DoTA, TADS, HTA y AI City Challenge). \NdatasetPropio{} estudios usan exclusivamente datos no p\'{u}blicos y \NdatosSimuladosAny{} usan datos simulados, en dos de ellos de forma exclusiva~\cite{chakraborty2023structural,park2025graph}. Ninguno de los datasets de sensores m\'{a}s usados (PeMS) incluye etiquetas de anomal\'{i}a verificadas, lo que explica el recurso a anomal\'{i}as inyectadas o a pseudo-etiquetas.

% Generado por suplementos/generar_conteos.py. No editar a mano.
\begin{table}[H]
\centering
\small
\caption{Tipo de dato y conjuntos de datos utilizados (RQ4).}\label{tab:datos}
\begin{tabularx}{\linewidth}{lcXX}
\toprule
\textbf{Tipo de dato} & \textbf{n} & \textbf{Estudios} & \textbf{Datasets} \\
\midrule
Series de sensores & 17 & S01, S02, S05, S08, S10, S12, S13, S16, S21, S31, S38, S41, S51, S63, S65, S69, S72 & 13 intersecciones, 66…, 45 incidentes reales, 4136…, 59 accidentes en el área de…, Dataset de simulación…, Dataset principal de…, Dataset real de la Agencia…, Dataset real de tráfico de…, Dataset sintético generado…, Detectores de lazo de…, Detectores de microondas de…, El artículo declara datos…, FT-AED, Freeway Traffic Anomalous…, I-880 traffic incident…, MnDOT / Regional…, Odense, Open ITS Open Traffic Data…, PeMS I-80, California, SUMO, Safety Pilot Model…, cada día contiene 480…, datos históricos de…, el número exacto de…, escenario sintético de red…, frecuencia de muestreo de…, más de 3.7 millones de…, no debe confundirse con el…, no se indican las fechas…, nombre público no reportado, país de la red simulada no…, series sintéticas derivadas…, velocidad y longitud). Los… \\
Movilidad agregada/OD & 2 & S30, S32 & Dataset real de monitoreo…, NYC Taxi Trips, publicado…, el artículo declara un…, se selecciona un subgrafo… \\
Trayectorias & 13 & S03, S11, S14, S15, S18, S20, S23, S24, S33, S34, S48, S54, S57 & 21,339 taxis, junio de 2013, Chengdu 2124763…, Chengdu, China…, Datasets de trayectorias…, Didi Chuxing privado, Dos conjuntos, Dos conjuntos reales, Dos conjuntos reales de…, ETC Shanghai, propio, Entrenamiento, MAAD, No especificado, Porto Taxi Trajectory, Porto se ajusta a la red…, Porto taxi, Tabla 4, p. 11), Taxis de Pekín, China, TransModeler, simulador…, Xi’an, 373054 trayectorias…, contiene 68.454.977 puntos…, el estudio emplea 1348…, entre 1000 y 2500…, frecuencia unificada de 10…, incluye 24 anomalías…, no se reportan cantidades…, proporcionados por Zhejiang…, prueba con 1795 vehículos y…, se utilizan datos simulados…, sección 5.2, p. 16. No se… \\
Video e im\'{a}genes & 29 & S04, S06, S07, S09, S17, S19, S25, S26, S27, S29, S35, S37, S42, S43, S44, S45, S46, S47, S49, S52, S55, S56, S60, S61, S64, S66, S67, S68, S71 & 127138 fotogramas normales…, 270 imágenes se destinan al…, 2782 videos reales de…, 357231 instancias obtenidas…, 488 fotogramas de accidente…, 4M-TAD, 594 imágenes recopiladas de…, 7160 fotogramas de…, 80 \% entrenamiento y 20 \%…, 80 secuencias de video con…, 94 720 fotogramas normales…, 94720 fotogramas normales…, CADP, Car Accident Detection and…, Car Crash Dataset, CarCrash, Caso 2, Conjunto de 3300 clips de…, Conjunto propio de 350…, Conjunto propio de 9000…, DAD, DADA-2000, Dataset UIT-ADrone…, Dataset propio de eventos…, Dataset propio de imágenes…, Dataset recopilado por los…, DoTA, DoTA, Detection of Traffic…, Fuentes declaradas, HWID12, IITH Road Accident Dataset, TAD, TRAF, conjunto de videos de…, Tabla 1, p. 5). Evaluación…, Traffic-Net, UA-DETRAC, UCF-Crime, conjunto público…, accidente 48 entrenamiento…, accidente situado…, actividad sospechosa 48…, cada ejemplo se representa…, cada video tiene 50…, conjuntos públicos de…, duración 5 s, 100…, grabaciones de un día…, imágenes web scraping, la referencia, los autores seleccionan y…, los autores trabajan con 11…, no especifican…, no se especifica su…, particiones de 2393…, particiones declaradas de…, procedencia geográfica no…, prueba 725, recopiladas de sitios web…, sección 2, p. 3. Tabla 1…, sección 3.1, pp. 70--71. Se…, sección 3.3, p. 70, y…, sección 4.2, pp. 9--10. No…, también informa 508…, ubicaciones geográficas no…, una clase de objeto… \\
Multimodal & 11 & S22, S28, S36, S39, S40, S50, S53, S58, S59, S62, S70 & 41630 corresponden al…, California PeMS, Estados…, Conjunto integrado de 12…, Cuatro conjuntos públicos…, Dataset propio de…, Dos conjuntos reales…, Driver Behavior and Route…, Driving Anomaly Dataset, Drone-Anomaly, 59 videos…, METR-LA y PEMS-BAY, TADS, UAH-DriveSet, dataset…, UCF-Crime, 1900 videos…, VTAD-highway, conjunto…, cada mapa resume en…, cantidad exacta de…, conjuntos públicos de…, desarrollo, 3 sesiones y…, el estudio utiliza 89…, entrenamiento con 15…, estas cantidades son…, junto con dos descripciones…, las unidades de carretera…, remite al trabajo original…, sección 4.1, pp. 4--5. Los…, sección 4.3, p. 9), sección 5.1, pp. 8--9. No se… \\
\bottomrule
\end{tabularx}
\end{table}


\subsection{RQ5: Limitaciones reportadas}

La limitaci\'{o}n m\'{a}s declarada es la generalizaci\'{o}n a otros escenarios (\Nlimgeneralizacion{} estudios), seguida del etiquetado de anomal\'{i}as (\Nlimetiquetado), el costo computacional y de despliegue (\Nlimcostodespliegue) y el uso de datos propios o de acceso limitado (\Nlimdatosreplicabilidad{}; Tabla~\ref{tab:limitaciones}). Llama la atenci\'{o}n lo que los autores casi no declaran: solo \Nlimmetricasnouniformes{} estudio menciona la falta de comparabilidad de las m\'{e}tricas y solo \Nliminterpretabilidad{} la interpretabilidad, aunque ambas aparecen como problemas en esta s\'{i}ntesis (Secci\'{o}n~\ref{sec:discusion}).

% Generado por suplementos/generar_conteos.py. No editar a mano.
\begingroup
\small
\begin{xltabular}{\linewidth}{lccX}
\caption{Limitaciones reportadas por los estudios primarios (RQ5).}\label{tab:limitaciones}\\
\toprule
\textbf{Limitaci\'{o}n declarada por los autores} & \textbf{n} & \textbf{\%} & \textbf{Estudios} \\
\midrule\endfirsthead
\toprule
\textbf{Limitaci\'{o}n declarada por los autores} & \textbf{n} & \textbf{\%} & \textbf{Estudios} \\
\midrule\endhead
\bottomrule\endfoot
Datos propios o de acceso limitado & 22 & 31~\% & S02, S05, S13, S14, S19, S20, S29, S30, S31, S32, S33, S40, S41, S44, S46, S54, S55, S56, S57, S59, S64, S69 \\
Etiquetado de anomal\'{i}as & 16 & 22~\% & S02, S06, S08, S11, S16, S17, S18, S30, S31, S34, S41, S50, S51, S55, S63, S71 \\
Generalizaci\'{o}n a otros escenarios & 34 & 47~\% & S02, S05, S13, S15, S18, S21, S23, S24, S25, S26, S27, S29, S36, S37, S41, S42, S45, S47, S49, S50, S52, S53, S56, S58, S60, S62, S63, S66, S67, S68, S69, S70, S71, S72 \\
M\'{e}tricas no comparables & 6 & 8~\% & S26, S29, S41, S49, S56, S66 \\
Costo computacional y despliegue & 29 & 40~\% & S02, S08, S09, S12, S17, S21, S22, S29, S35, S36, S37, S38, S40, S41, S43, S44, S45, S46, S47, S53, S57, S60, S61, S62, S64, S66, S68, S70, S72 \\
Interpretabilidad & 4 & 6~\% & S22, S32, S62, S67 \\
Sin limitaciones declaradas & 9 & 12~\% & S01, S03, S04, S07, S10, S28, S39, S48, S65 \\
\end{xltabular}
\vspace{-6pt}{\footnotesize\raggedright Un estudio puede declarar varias limitaciones.\par}
\endgroup


\subsection{RQ6: L\'{i}neas base y ventaja reportada}
\label{sec:rq6}

Todos los estudios comparan su propuesta con alguna l\'{i}nea base, pero con alcance desigual (Tabla~\ref{tab:baselines}). \NbaselineNoDL{} estudios incluyen al menos una l\'{i}nea base que no es DL (estad\'{i}stica, de aprendizaje autom\'{a}tico convencional o basada en reglas), lo que permite valorar el aporte del DL frente a alternativas m\'{a}s simples. En otros casos, la comparaci\'{o}n se limita a ablaciones del propio modelo~\cite{huang2020intelligent,khan2022anomaly}, a cifras copiadas de art\'{i}culos con otros datos~\cite{jadhav2026multistage} o a un componente auxiliar y no a la detecci\'{o}n~\cite{zhao2024multiadjacent}. Las ventajas reportadas frente a la mejor l\'{i}nea base suelen ser moderadas (de 1 a 6 puntos porcentuales en la m\'{e}trica principal) cuando la comparaci\'{o}n es con modelos recientes bajo el mismo protocolo~\cite{chai2026road,park2025graph,abdallah2025aipurban,duan2025multimodal}, y mayores en algunos casos, como +0,15 de AUC frente a redes de grafos gen\'{e}ricas~\cite{hu2022detecting} o $\approx$\,+17 puntos de F1 frente a m\'{e}todos de umbral~\cite{cai2021road}. Hay resultados desfavorables que los estudios reportan, pero no destacan: con ventanas cortas, Autoformer supera al modelo propuesto por Park et al.~\cite{park2025graph}, y el FPR de Chakraborty et al.~\cite{chakraborty2023structural} es peor que el de su mejor l\'{i}nea base. Fuera del corpus, Taheri Sarteshnizi et al.~\cite{taheri2026applicability} eval\'{u}an detectores de DL sobre datos reales de volumen vehicular con etiquetas expertas y encuentran que un m\'{e}todo simple basado en distancias supera a todos ellos (F1 = 0,676 frente a 0,278 del mejor detector profundo). Este estudio se excluy\'{o} por CE-4, pero es relevante para valorar la ventaja del DL.

\begin{landscape}
\footnotesize
\begin{longtable}{lP{7.5cm}cP{8.5cm}}
\caption{L\'{i}neas base y ventaja frente a la mejor de ellas (RQ6).\label{tab:baselines}}\\
\toprule
\textbf{ID} & \textbf{L\'{i}neas base} & \textbf{No DL} & \textbf{Ventaja frente a la mejor l\'{i}nea base} \\
\midrule\endfirsthead
\toprule
\textbf{ID} & \textbf{L\'{i}neas base} & \textbf{No DL} & \textbf{Ventaja frente a la mejor l\'{i}nea base} \\
\midrule\endhead
\bottomrule\endfoot
\tabinput{generado/baselines.tex}
\end{longtable}
\end{landscape}

\paragraph{An\'{a}lisis de sensibilidad por calidad.} Restringida a los \NaltaCalidad{} estudios de calidad alta (\IDaltaCalidad), la s\'{i}ntesis mantiene sus conclusiones principales: \NaltaHibrido{} de ellos son h\'{i}bridos y la arquitectura sigue el tipo de dato. Sin embargo, solo \NaltaEtiqReal{} eval\'{u}an con etiquetas reales o anotadas, de modo que la escasez de etiquetas verificadas afecta tambi\'{e}n a los mejores estudios del corpus.

%=============================================================
\section{Discusi\'{o}n}
\label{sec:discusion}
%=============================================================

\subsection{Interpretaci\'{o}n de los hallazgos y contraste con revisiones previas}

\paragraph{El dato antes que el modelo.} El hallazgo principal es que la familia de arquitectura se explica mejor por el tipo de dato que por el tipo de anomal\'{i}a. Las revisiones previas trataban por separado el video~\cite{santhosh2020anomaly,fang2024vision}, las trayectorias~\cite{belhadi2020trajectory} y las series de sensores~\cite{djenouri2019survey}, por lo que no pod\'{i}an observar esta relaci\'{o}n. Al integrarlos, esta revisi\'{o}n muestra que no existe una arquitectura dominante para el problema en general: las redes convolucionales son la opci\'{o}n natural cuando la anomal\'{i}a es visible en la imagen, las de grafos cuando el dato tiene estructura de red y los autoencoders cuando faltan etiquetas. Este resultado coincide con la observaci\'{o}n de Darban et al.~\cite{darban2024deep} para series temporales gen\'{e}ricas: la idoneidad de un detector depende de la estructura de los datos y del tipo de anomal\'{i}a, no de su complejidad.

\paragraph{Hibridaci\'{o}n y atenci\'{o}n.} La hibridaci\'{o}n es la norma y no la excepci\'{o}n, y la atenci\'{o}n o los Transformers se han generalizado desde 2024. Djenouri et al.~\cite{djenouri2019survey} y Zhang et al.~\cite{zhang2020urban} reportaban un campo dominado por m\'{e}todos estad\'{i}sticos y de miner\'{i}a de datos, por lo que el cambio hacia modelos profundos h\'{i}bridos es claro. Sin embargo, la evidencia de que esta complejidad mejore la detecci\'{o}n en condiciones reales es limitada: las ventajas frente a l\'{i}neas base recientes son moderadas, y Taheri Sarteshnizi et al.~\cite{taheri2026applicability} encuentran que, con etiquetas reales, m\'{e}todos simples pueden superar a los profundos.

\paragraph{Etiquetas y protocolo.} La aportaci\'{o}n m\'{a}s espec\'{i}fica de esta revisi\'{o}n es mostrar que la comparabilidad depende de la procedencia de las etiquetas y del protocolo. Ninguna de las revisiones previas analiza estos aspectos de forma estructurada. Que un mismo dataset (PeMS) se use con cuatro definiciones distintas de anomal\'{i}a, y que con anomal\'{i}as inyectadas se obtengan AUC perfectos, indica que las cifras de desempe\~{n}o publicadas no pueden ordenarse en un ranking. Esto coincide con la cr\'{i}tica de Darban et al.~\cite{darban2024deep} sobre los benchmarks de series temporales y con las recomendaciones REFORMS para estudios de aprendizaje autom\'{a}tico~\cite{kapoor2024reforms}.

\paragraph{Alcance de las conclusiones.} Los estudios incluidos aplican DL con resultados favorables en escenarios controlados, pero con validaci\'{o}n limitada en contextos reales: \NdatosSimuladosAny{} usan datos simulados, \NetiqNoReales{} no eval\'{u}an con etiquetas reales y solo \QASeissi{} validan en m\'{a}s de un escenario independiente. Por tanto, no se puede afirmar que el DL haya alcanzado madurez operativa para la detecci\'{o}n de anomal\'{i}as de tr\'{a}fico.

\subsection{Implicaciones para sistemas inteligentes de transporte}

Para un operador de tr\'{a}fico, la elecci\'{o}n del modelo deber\'{i}a partir del dato disponible: video en intersecciones y t\'{u}neles, detectores en autopistas, trayectorias de flotas en redes urbanas. Antes de adoptar un modelo publicado, conviene verificar con qu\'{e} etiquetas se evalu\'{o}, con qu\'{e} proporci\'{o}n de anomal\'{i}as y con qu\'{e} unidad de evaluaci\'{o}n, porque un F1 obtenido con anomal\'{i}as inyectadas no anticipa el desempe\~{n}o con incidentes reales. Las m\'{e}tricas operativas (latencia, tiempo medio de detecci\'{o}n, tasa de falsas alarmas por hora) se reportan poco y son las que determinan el valor pr\'{a}ctico de una alerta.

\subsection{Brechas y agenda de investigaci\'{o}n}

\begin{enumerate}[label=(\arabic*)]
  \item \textbf{Benchmarks con etiquetas verificadas.} Solo \NetiqVerificadas{} estudios eval\'{u}an con etiquetas reales verificadas. Se necesitan conjuntos de datos p\'{u}blicos de sensores y trayectorias con incidentes registrados por operadores, en la l\'{i}nea de AnoLT~\cite{taheri2026applicability}.
  \item \textbf{Reporte m\'{i}nimo del protocolo.} Se propone reportar siempre la procedencia de las etiquetas, la proporci\'{o}n de anomal\'{i}as, la unidad de evaluaci\'{o}n, el criterio de umbral, la partici\'{o}n temporal y la variabilidad entre ejecuciones, junto con PR-AUC cuando las clases est\'{e}n desbalanceadas.
  \item \textbf{Validaci\'{o}n externa.} Pocos estudios eval\'{u}an en m\'{a}s de una ciudad o v\'{i}a. La transferencia entre ciudades y la adaptaci\'{o}n de dominio siguen sin evaluarse sistem\'{a}ticamente.
  \item \textbf{Reproducibilidad.} Solo \NcodigoURL{} estudios publican su c\'{o}digo. La publicaci\'{o}n de c\'{o}digo y datos deber\'{i}a acompa\~{n}ar a los resultados.
  \item \textbf{Interpretabilidad.} Casi ning\'{u}n estudio la aborda ni la declara como limitaci\'{o}n, pese a que los operadores necesitan entender el origen de una alerta.
\end{enumerate}

\subsection{Amenazas a la validez}
\label{sec:amenazas}

\paragraph{Validez de construcci\'{o}n (b\'{u}squeda).} La cadena no fue validada por un especialista en informaci\'{o}n externo. Para mitigarlo, se valid\'{o} con un conjunto de control (\NsensibOK{} de \Nestudios{} estudios recuperados por la cadena definitiva). Dos estudios solo se recuperan con b\'{u}squeda en texto completo, lo que indica que pueden existir otros estudios relevantes que no mencionan DL en el resumen. La revisi\'{o}n se limita a cinco bases y no incluye literatura gris ni \textit{preprints}, lo que excluye trabajos recientes a\'{u}n no publicados. Adem\'{a}s, solo se incluyen estudios de acceso abierto: esto garantiza que cualquier lector pueda verificar la extracci\'{o}n, pero deja fuera estudios relevantes publicados bajo suscripci\'{o}n, y la proporci\'{o}n de acceso abierto var\'{i}a entre editoriales y pa\'{i}ses, por lo que el corpus puede no representar por igual a toda la literatura.

\paragraph{Validez interna (selecci\'{o}n y extracci\'{o}n).} La versi\'{o}n anterior de la revisi\'{o}n aplic\'{o} los criterios de forma no uniforme: incluy\'{o} un \textit{preprint} como excepci\'{o}n y estudios sin evaluaci\'{o}n cuantitativa de la detecci\'{o}n. En esta versi\'{o}n los criterios se aplicaron a todos los estudios y los casos l\'{i}mite se documentan en el Anexo~\ref{ap:limite}. La principal amenaza de esta versi\'{o}n es que la selecci\'{o}n, la extracci\'{o}n y la evaluaci\'{o}n de calidad las realiz\'{o} un solo revisor, lo que PRISMA y AMSTAR~2 consideran un riesgo de sesgo de selecci\'{o}n. Se mitig\'{o} con criterios operativos aplicados en orden fijo, exclusiones a texto completo respaldadas por citas, un asistente de cribado basado en reglas cuyo acuerdo con el revisor se reporta, y un re-cribado ciego para medir la consistencia intra-revisor (Secci\'{o}n~\ref{sec:metodo}). Estas medidas no sustituyen a un segundo revisor independiente. La extracci\'{o}n se verific\'{o} contra el texto completo, y durante esa verificaci\'{o}n se encontraron errores en las fichas iniciales (fichas duplicadas, categor\'{i}as y datasets mal asignados), que se corrigieron.

\paragraph{Validez externa.} El corpus es peque\~{n}o (\Nestudios{} estudios), est\'{a} concentrado en datos de China y Estados Unidos y se limita a estudios en ingl\'{e}s. El a\~{n}o 2026 es parcial. Las proporciones deben interpretarse como descripci\'{o}n del corpus y no como estimaciones de la literatura completa.

\paragraph{Validez de conclusi\'{o}n.} La heterogeneidad de tareas, datos, etiquetas y protocolos impide el metaan\'{a}lisis y la valoraci\'{o}n formal de la certeza de la evidencia. Por ello, las conclusiones sobre la ventaja del DL se formulan con prudencia. Varios estudios presentan inconsistencias internas entre texto y tablas, registradas en la matriz de extracci\'{o}n, que reducen la confianza en sus resultados num\'{e}ricos.

%=============================================================
\section{Conclusiones}
\label{sec:conclusiones}
%=============================================================

Esta revisi\'{o}n sistem\'{a}tica analiz\'{o} \Nestudios{} estudios (\AnioMin--\AnioMax) que aplican \textit{Deep Learning} a la detecci\'{o}n de anomal\'{i}as en tr\'{a}fico vehicular urbano y de autopista a partir de sensores, trayectorias, movilidad agregada y video.

\textbf{RQ1.} Predominan las redes recurrentes, convolucionales y de atenci\'{o}n, casi siempre combinadas en modelos h\'{i}bridos. La familia de arquitectura sigue al tipo de dato: convolucionales en video, grafos en datos con estructura de red y autoencoders cuando faltan etiquetas. Desde 2024 la atenci\'{o}n y los Transformers est\'{a}n presentes en la mayor\'{i}a de los estudios.

\textbf{RQ2.} Las m\'{e}tricas m\'{a}s usadas son \textit{accuracy}, recall, AUC-ROC, precision y F1, con definiciones a veces no est\'{a}ndar. El principal obst\'{a}culo para comparar estudios es el protocolo: la procedencia de las etiquetas (reales, inyectadas o pseudo-etiquetas), la proporci\'{o}n de anomal\'{i}as y el criterio de umbral var\'{i}an entre estudios y a menudo no se reportan.

\textbf{RQ3.} Las anomal\'{i}as m\'{a}s estudiadas son los accidentes e incidentes (\PcatA), seguidos de la congesti\'{o}n o flujo an\'{o}malo, las trayectorias peligrosas y las disrupciones.

\textbf{RQ4.} Los datasets m\'{a}s reutilizados (PeMS y taxis) carecen de etiquetas de anomal\'{i}a verificadas. En video predominan conjuntos espec\'{i}ficos de cada estudio, y una parte del corpus depende de datos simulados o no p\'{u}blicos.

\textbf{RQ5.} Los autores declaran sobre todo problemas de generalizaci\'{o}n, etiquetado y costo computacional. Casi ninguno declara la falta de comparabilidad de las m\'{e}tricas o la interpretabilidad.

\textbf{RQ6.} Todos los estudios comparan con l\'{i}neas base, pero solo \NbaselineNoDL{} de \Nestudios{} incluyen una alternativa no profunda y varios se limitan a ablaciones o a cifras de otros art\'{i}culos. Las ventajas frente a modelos recientes son moderadas.

Como implicaci\'{o}n pr\'{a}ctica, la elecci\'{o}n del modelo deber\'{i}a partir del dato disponible y de c\'{o}mo se evalu\'{o} el modelo, no solo de la cifra de desempe\~{n}o reportada. Como trabajo futuro se proponen benchmarks p\'{u}blicos con etiquetas verificadas, un reporte m\'{i}nimo del protocolo de evaluaci\'{o}n y la validaci\'{o}n entre ciudades.

%=============================================================
\section*{Declaraciones}
%=============================================================

\textbf{Registro y protocolo:} protocolo registrado en \RegistroOSF. La lista de verificaci\'{o}n PRISMA 2020 se presenta en el Anexo~\ref{ap:prisma}.

\textbf{Disponibilidad de datos y c\'{o}digo:} la matriz de extracci\'{o}n, la evaluaci\'{o}n de calidad por \'{i}tem y estudio, las cadenas de b\'{u}squeda, la planilla de cribado y los scripts que generan todas las cifras (\texttt{generar\_conteos.py} y \texttt{prueba\_sensibilidad\_cadena.py}) est\'{a}n disponibles en \RepositorioDatos.

\textbf{Conflicto de intereses:} el autor declara no tener conflicto de intereses.

\textbf{Financiamiento:} esta investigaci\'{o}n no recibi\'{o} financiamiento externo.

\textbf{Uso de inteligencia artificial generativa:} el autor utiliz\'{o} Claude (Anthropic) como asistente para verificar la extracci\'{o}n de datos contra los textos completos, proponer una decisi\'{o}n justificada para cada registro en el cribado por t\'{i}tulo y resumen, proponer los puntajes de calidad con su evidencia, elaborar los scripts de deduplicaci\'{o}n, cribado asistido y conteo, y revisar la redacci\'{o}n. Todas las decisiones de inclusi\'{o}n, los puntajes y el texto fueron revisados y validados por el autor, que asume la responsabilidad del trabajo.

\textbf{Contribuciones (CRediT):} \textbf{Mayta Arotaype, W. A.}: conceptualizaci\'{o}n, metodolog\'{i}a, b\'{u}squeda, selecci\'{o}n de estudios, curaci\'{o}n de datos, an\'{a}lisis formal, visualizaci\'{o}n y redacci\'{o}n.

%=============================================================
\bibliographystyle{IEEEtran}
\bibliography{referencias}
%=============================================================

\appendix

\section{Protocolo de revisi\'{o}n}
\label{ap:protocolo}

\subsection*{Cadena de b\'{u}squeda por base de datos}

Las cadenas se reproducen literalmente. Los filtros de a\~{n}o, tipo de documento e idioma se aplicaron con las opciones de la interfaz en Web of Science y ScienceDirect, y con un script sobre los metadatos exportados en Scopus, IEEE Xplore y ACM (Tabla~\ref{tab:bases}).

{\footnotesize
\noindent\textbf{Scopus} (b\'{u}squeda avanzada)
\begin{verbatim}
TITLE-ABS-KEY ( ( anomal* OR abnormal* OR "outlier detection" OR "incident
 detection" OR "accident detection" OR "accident prediction" OR "accident
 anticipation" OR "road closure*" OR "dangerous driving" OR "near-accident*"
 OR "near-miss*" OR "extreme event*" ) AND ( "deep learning" OR "deep neural"
 OR "neural network*" OR lstm OR gru OR cnn OR convolutional OR autoencoder*
 OR transformer* OR "graph neural" OR "self-attention" OR "attention
 mechanism" ) AND ( "road traffic" OR "urban traffic" OR "vehicular traffic"
 OR "highway traffic" OR "freeway traffic" OR "expressway traffic" OR "traffic
 flow*" OR "traffic speed" OR "traffic volume*" OR "traffic data" OR "traffic
 monitoring" OR "traffic video*" OR "traffic surveillance" OR "traffic
 camera*" OR "traffic scene*" OR "traffic scenario*" OR "traffic event*" OR
 "traffic incident*" OR "traffic accident*" OR "traffic anomal*" OR "traffic
 congestion" OR "road accident*" OR "road closure*" OR "intelligent
 transportation" OR "transportation network*" OR "road network*" OR "vehicle
 trajector*" OR "vehicle behavio*" OR "driving behavio*" OR taxi* ) ) AND
 PUBYEAR > 2019 AND PUBYEAR < 2027 AND ( LIMIT-TO ( DOCTYPE , "ar" ) OR
 LIMIT-TO ( DOCTYPE , "cp" ) ) AND ( LIMIT-TO ( LANGUAGE , "English" ) OR
 LIMIT-TO ( LANGUAGE , "Spanish" ) )
\end{verbatim}

\noindent\textbf{Web of Science Core Collection} (TS = t\'{i}tulo, resumen, palabras clave y Keywords Plus)
\begin{verbatim}
TS=((anomal* OR abnormal* OR "outlier detection" OR "incident detection" OR
 "accident detection" OR "accident prediction" OR "accident anticipation" OR
 "road closure*" OR "dangerous driving" OR "near-accident*" OR "near-miss*" OR
 "extreme event*") AND ("deep learning" OR "deep neural" OR "neural network*"
 OR LSTM OR GRU OR CNN OR convolutional OR autoencoder* OR transformer* OR
 "graph neural" OR "self-attention" OR "attention mechanism") AND ("road
 traffic" OR "urban traffic" OR "vehicular traffic" OR "highway traffic" OR
 "freeway traffic" OR "expressway traffic" OR "traffic flow*" OR "traffic
 speed" OR "traffic volume*" OR "traffic data" OR "traffic monitoring" OR
 "traffic video*" OR "traffic surveillance" OR "traffic camera*" OR "traffic
 scene*" OR "traffic scenario*" OR "traffic event*" OR "traffic incident*" OR
 "traffic accident*" OR "traffic anomal*" OR "traffic congestion" OR "road
 accident*" OR "road closure*" OR "intelligent transportation" OR
 "transportation network*" OR "road network*" OR "vehicle trajector*" OR
 "vehicle behavio*" OR "driving behavio*" OR taxi*))
Filtros (panel izquierdo): Publication Years 2020-2026; Document Types:
 Article, Proceeding Paper; Languages: English, Spanish.
\end{verbatim}

\noindent\textbf{IEEE Xplore} (Command Search, \textit{All Metadata}; m\'{a}ximo 25 t\'{e}rminos por consulta)
\begin{verbatim}
Consulta 1: ("All Metadata":anomal* OR "All Metadata":abnormal* OR "All
 Metadata":"incident detection" OR "All Metadata":"accident detection" OR "All
 Metadata":"accident prediction" OR "All Metadata":"road closure" OR "All
 Metadata":"dangerous driving" OR "All Metadata":"near-miss" OR "All
 Metadata":"extreme event") AND ("All Metadata":"deep learning" OR "All
 Metadata":"neural network" OR "All Metadata":LSTM OR "All Metadata":CNN OR
 "All Metadata":convolutional OR "All Metadata":autoencoder OR "All
 Metadata":transformer OR "All Metadata":"graph neural") AND ("All
 Metadata":"road traffic" OR "All Metadata":"urban traffic" OR "All
 Metadata":"highway traffic" OR "All Metadata":"traffic flow" OR "All
 Metadata":"traffic video" OR "All Metadata":"traffic surveillance" OR "All
 Metadata":"traffic incident" OR "All Metadata":"traffic accident")
Consulta 2: ("All Metadata":anomal* OR "All Metadata":abnormal* OR "All
 Metadata":"incident detection" OR "All Metadata":"accident detection" OR "All
 Metadata":"accident prediction" OR "All Metadata":"road closure" OR "All
 Metadata":"dangerous driving" OR "All Metadata":"near-miss" OR "All
 Metadata":"extreme event") AND ("All Metadata":"deep learning" OR "All
 Metadata":"neural network" OR "All Metadata":LSTM OR "All Metadata":CNN OR
 "All Metadata":convolutional OR "All Metadata":autoencoder OR "All
 Metadata":transformer OR "All Metadata":"graph neural") AND ("All
 Metadata":"traffic anomaly" OR "All Metadata":"traffic congestion" OR "All
 Metadata":"traffic data" OR "All Metadata":"traffic monitoring" OR "All
 Metadata":"intelligent transportation" OR "All Metadata":"road network" OR
 "All Metadata":"vehicle trajectory" OR "All Metadata":"driving behavior")
Consulta 3: ("All Metadata":anomal* OR "All Metadata":abnormal* OR "All
 Metadata":"incident detection" OR "All Metadata":"accident detection" OR "All
 Metadata":"accident prediction" OR "All Metadata":"road closure" OR "All
 Metadata":"dangerous driving" OR "All Metadata":"near-miss" OR "All
 Metadata":"extreme event") AND ("All Metadata":"deep learning" OR "All
 Metadata":"neural network" OR "All Metadata":LSTM OR "All Metadata":CNN OR
 "All Metadata":convolutional OR "All Metadata":autoencoder OR "All
 Metadata":transformer OR "All Metadata":"graph neural") AND ("All
 Metadata":"vehicular traffic" OR "All Metadata":"freeway traffic" OR "All
 Metadata":"traffic speed" OR "All Metadata":"traffic volume" OR "All
 Metadata":"traffic scene" OR "All Metadata":"traffic event" OR "All
 Metadata":"road accident" OR "All Metadata":taxi)
Filtros: Year 2020-2026; Content Type: Journals, Conferences.
\end{verbatim}

\noindent\textbf{ScienceDirect} (t\'{i}tulo, resumen y palabras clave; m\'{a}ximo 8 operadores, sin comodines)
\begin{verbatim}
Consulta 1: ("anomaly detection" OR "incident detection" OR "accident
 detection") AND ("deep learning" OR "neural network") AND ("road traffic" OR
 "urban traffic" OR "traffic flow")
Consulta 2: ("anomaly detection" OR "incident detection" OR "accident
 detection") AND ("deep learning" OR "neural network") AND ("traffic video" OR
 "traffic surveillance" OR "traffic accident")
Consulta 3: ("anomaly detection" OR "incident detection" OR "accident
 detection") AND ("deep learning" OR "neural network") AND ("traffic incident"
 OR "traffic congestion" OR "intelligent transportation")
Consulta 4: ("anomaly detection" OR "incident detection" OR "accident
 detection") AND ("deep learning" OR "neural network") AND ("vehicle
 trajectory" OR "road network" OR "highway traffic")
Consulta 5: ("abnormal" OR "anomalous") AND ("deep learning" OR "neural
 network") AND ("road traffic" OR "traffic flow" OR "vehicle trajectory")
Filtros: Years 2020-2026; Article type: Research articles.
\end{verbatim}

\noindent\textbf{ACM Digital Library} (campo Abstract)
\begin{verbatim}
[[Abstract: anomal*] OR [Abstract: abnormal*] OR [Abstract: "outlier
 detection"] OR [Abstract: "incident detection"] OR [Abstract: "accident
 detection"] OR [Abstract: "accident prediction"] OR [Abstract: "road
 closure*"] OR [Abstract: "dangerous driving"] OR [Abstract: "near-accident*"]
 OR [Abstract: "near-miss*"] OR [Abstract: "extreme event*"]] AND [[Abstract:
 "deep learning"] OR [Abstract: "deep neural"] OR [Abstract: "neural
 network*"] OR [Abstract: lstm] OR [Abstract: gru] OR [Abstract: cnn] OR
 [Abstract: convolutional] OR [Abstract: autoencoder*] OR [Abstract:
 transformer*] OR [Abstract: "graph neural"] OR [Abstract: "attention
 mechanism"]] AND [[Abstract: "road traffic"] OR [Abstract: "urban traffic"]
 OR [Abstract: "vehicular traffic"] OR [Abstract: "highway traffic"] OR
 [Abstract: "freeway traffic"] OR [Abstract: "traffic flow*"] OR [Abstract:
 "traffic speed"] OR [Abstract: "traffic volume*"] OR [Abstract: "traffic
 data"] OR [Abstract: "traffic monitoring"] OR [Abstract: "traffic video*"] OR
 [Abstract: "traffic surveillance"] OR [Abstract: "traffic camera*"] OR
 [Abstract: "traffic scene*"] OR [Abstract: "traffic event*"] OR [Abstract:
 "traffic incident*"] OR [Abstract: "traffic accident*"] OR [Abstract:
 "traffic anomal*"] OR [Abstract: "traffic congestion"] OR [Abstract: "road
 accident*"] OR [Abstract: "intelligent transportation"] OR [Abstract: "road
 network*"] OR [Abstract: "vehicle trajector*"] OR [Abstract: "driving
 behavio*"] OR [Abstract: taxi*]] AND [E-Publication Date: (01/01/2020 TO
 12/31/2026)]
Filtros: Content type: Research Article.
\end{verbatim}
}

\subsection*{Historial del protocolo}

\begin{table}[H]
\centering
\small
\caption{Historial de versiones del protocolo.}
\label{tab:historial}
\begin{tabularx}{\linewidth}{llY}
\toprule
\textbf{Versi\'{o}n} & \textbf{Fecha} & \textbf{Cambio} \\
\midrule
V1.0 & \pend{fecha} & Definici\'{o}n del tema, preguntas de investigaci\'{o}n y PICOC. \\
V1.1 & \pend{fecha} & Cadenas por base (seis fuentes, incluidas MDPI y Taylor \& Francis). \\
V1.2 & \pend{fecha} & Cribado y eliminaci\'{o}n de duplicados en Parsif.al. \\
V1.3 & \pend{fecha} & Evaluaci\'{o}n de calidad (instrumento de 6 \'{i}tems) y extracci\'{o}n; 25 estudios. \\
V2.1 & \pend{fecha} & Bloque C de la cadena acotado a frases ancladas al tr\'{a}fico vial (se retiran \textit{vehicle*}, \textit{highway*}, \textit{freeway*}, \textit{expressway*} y \textit{outlier*} sueltos) para reducir el ruido; la recuperaci\'{o}n del conjunto de control se mantiene (19 de 21). \\
V2.0 & \pend{fecha} & Respuesta a la revisi\'{o}n: alcance ampliado a video y autopistas; cinco bases (se a\~{n}ade IEEE Xplore y se retiran MDPI y Taylor \& Francis por ser editoriales); cadena validada con un conjunto de control; regla general para \textit{preprints} (CE-5); CI-5 y CI-6 precisados; nuevo instrumento de calidad de 8 \'{i}tems; nueva RQ6; cribado por un revisor con asistente basado en reglas y re-cribado ciego. \\
\bottomrule
\end{tabularx}
\end{table}

\section{Estudios incluidos}
\label{ap:estudios}

\begin{landscape}
\footnotesize
\begin{longtable}{lllP{2.8cm}cP{3.0cm}P{5.5cm}P{3.0cm}}
\caption{Estudios incluidos con su identificador en la versi\'{o}n anterior, DOI y base de recuperaci\'{o}n.\label{tab:ap_estudios}}\\
\toprule
\textbf{ID} & \textbf{ID v1} & \textbf{Ref.} & \textbf{Autores} & \textbf{A\~{n}o} & \textbf{Editorial} & \textbf{DOI} & \textbf{Base de recuperaci\'{o}n} \\
\midrule\endfirsthead
\toprule
\textbf{ID} & \textbf{ID v1} & \textbf{Ref.} & \textbf{Autores} & \textbf{A\~{n}o} & \textbf{Editorial} & \textbf{DOI} & \textbf{Base de recuperaci\'{o}n} \\
\midrule\endhead
\bottomrule\endfoot
\tabinput{generado/lista_estudios.tex}
\end{longtable}
\end{landscape}

\section{Decisiones sobre casos l\'{i}mite}
\label{ap:limite}

\begin{small}
\begin{longtable}{P{2.6cm}P{2.2cm}P{9.6cm}}
\caption{Decisiones de elegibilidad sobre estudios dudosos, con el criterio aplicado.\label{tab:limite}}\\
\toprule
\textbf{Estudio} & \textbf{Decisi\'{o}n} & \textbf{Justificaci\'{o}n} \\
\midrule\endfirsthead
\toprule
\textbf{Estudio} & \textbf{Decisi\'{o}n} & \textbf{Justificaci\'{o}n} \\
\midrule\endhead
\bottomrule\endfoot
\tabinput{generado/limite.tex}
\end{longtable}
\end{small}

\section{Lista de verificaci\'{o}n PRISMA 2020}
\label{ap:prisma}

\begin{small}
\begin{longtable}{lP{8.8cm}P{4.0cm}}
\caption{Ubicaci\'{o}n de los \'{i}tems de PRISMA 2020 en el manuscrito.\label{tab:prisma_check}}\\
\toprule
\textbf{\'{I}tem} & \textbf{Requisito} & \textbf{Ubicaci\'{o}n} \\
\midrule\endfirsthead
\toprule
\textbf{\'{I}tem} & \textbf{Requisito} & \textbf{Ubicaci\'{o}n} \\
\midrule\endhead
\bottomrule\endfoot
1 & T\'{i}tulo que identifica el trabajo como revisi\'{o}n sistem\'{a}tica & T\'{i}tulo \\
2 & Resumen estructurado (PRISMA for Abstracts) & Resumen \\
3 & Justificaci\'{o}n en relaci\'{o}n con el conocimiento existente & Sec.~\ref{sec:intro}, \ref{sec:previas}; Tabla~\ref{tab:previas} \\
4 & Objetivos o preguntas & Sec.~\ref{sec:intro} (RQ1--RQ6) \\
5 & Criterios de elegibilidad & Tabla~\ref{tab:criterios} \\
6 & Fuentes de informaci\'{o}n y fecha de \'{u}ltima b\'{u}squeda & Sec.~\ref{sec:busqueda}; Tabla~\ref{tab:bases} \\
7 & Estrategia de b\'{u}squeda completa & Anexo~\ref{ap:protocolo} \\
8 & Proceso de selecci\'{o}n & Sec.~\ref{sec:metodo} (Proceso de selecci\'{o}n) \\
9 & Proceso de extracci\'{o}n & Sec.~\ref{sec:metodo} (Extracci\'{o}n de datos) \\
10 & Variables extra\'{i}das & Tabla~\ref{tab:formulario} \\
11 & Evaluaci\'{o}n del riesgo de sesgo & Sec.~\ref{sec:metodo}; Tabla~\ref{tab:cqi} \\
12 & Medidas de efecto & No aplica (s\'{i}ntesis sin metaan\'{a}lisis); Sec.~\ref{sec:metodo} (S\'{i}ntesis) \\
13a--f & M\'{e}todos de s\'{i}ntesis & Sec.~\ref{sec:metodo} (S\'{i}ntesis) \\
14 & Sesgo de reporte & Sec.~\ref{sec:metodo} (S\'{i}ntesis), justificaci\'{o}n \\
15 & Certeza de la evidencia & Sec.~\ref{sec:metodo} (S\'{i}ntesis), justificaci\'{o}n \\
16a & Resultados de la b\'{u}squeda y selecci\'{o}n & Figura~\ref{fig:prisma} \\
16b & Estudios excluidos que parec\'{i}an cumplir los criterios & Anexo~\ref{ap:limite} \\
17 & Caracter\'{i}sticas de los estudios & Tabla~\ref{tab:caracteristicas} \\
18 & Riesgo de sesgo por estudio & Tabla~\ref{tab:calidad} \\
19 & Resultados de cada estudio & Tablas~\ref{tab:caracteristicas} y \ref{tab:baselines} \\
20a--d & Resultados de las s\'{i}ntesis & Sec.~\ref{sec:resultados} (RQ1--RQ6) \\
21 & Sesgo de reporte en la s\'{i}ntesis & Sec.~\ref{sec:rq6} \\
22 & Certeza de la evidencia & Sec.~\ref{sec:amenazas} \\
23a--d & Discusi\'{o}n, limitaciones e implicaciones & Sec.~\ref{sec:discusion} \\
24a--c & Registro y protocolo & Declaraciones; Anexo~\ref{ap:protocolo} \\
25 & Financiamiento & Declaraciones \\
26 & Conflicto de intereses & Declaraciones \\
27 & Disponibilidad de datos, c\'{o}digo y materiales & Declaraciones \\
\end{longtable}
\end{small}

\end{document}
